\providecommand{\zerJ}{\mathcal{J}}
\providecommand{\zerK}{\mathcal{K}}
\providecommand{\zersinc}{\operatorname{sinc}}
\providecommand{\zergE}{\gamma_{\mathrm{E}}}
\providecommand{\zerLL}{\mathcal{L}}
\providecommand{\zerconv}{\operatorname{conv}}
\providecommand{\zerdiag}{\operatorname{diag}}
\providecommand{\zerSym}{\mathfrak{D}}
\providecommand{\zerLog}{\operatorname{Log}}
\section{Degree-two targets: Artin, Hecke and modular \texorpdfstring{$L$}{L}-functions}\label{sec:hec}

\providecommand{\hecsgn}{\operatorname{sgn}}
\providecommand{\hecVar}{\operatorname{Var}}
\providecommand{\hecCov}{\operatorname{Cov}}
\providecommand{\hecE}{\mathbb{E}}
\providecommand{\hecind}{\mathbf{1}}
% Local line-breaking tolerance for this section; restored at its end.
\edef\hecsavedemergencystretch{\the\emergencystretch}
\emergencystretch=2em\relax

The input theorem, \cref{src:thm:dgrh}, concerns $\zeta$ and the Dirichlet $L$-functions. The simplest
$L$-functions it does not reach have degree two over $\Q$: Hecke $L$-functions of ideal class characters
of imaginary quadratic fields (equivalently, two-dimensional dihedral Artin $L$-functions), Hecke
$L$-functions of characters of infinite order of $\Q(i)$, and the $L$-functions of holomorphic
modular forms together with their symmetric powers. This section collects what the one-sided method
gives for these targets. For each of them we state the exact prime estimate that is equivalent to
its Riemann hypothesis, we solve the extremal problems that govern the admissible prime weights, and we
record the unconditional results obtained along the way: real zeros, Jensen polynomials, oscillation of
the prime sums, and obstructions to realizing the weights by congruence conditions. Every statement that
mentions a Riemann hypothesis outside the scope of \cref{src:thm:dgrh} is a reduction; the prime
estimates themselves are open. The final subsection treats two questions about Dirichlet characters
that are settled, or sharply delimited, by \cref{src:thm:dgrh}.

\subsection{Conventions and the Landau step}\label{hec:ss:landau}
{\emergencystretch=\textwidth\sources{grh_transfer_analysis; grh_beyond_dirichlet_2026-11-25; grh_dihedral_47_71; grh_method_transfers_sharp_torsion}}

For an $L$-function $L(s)$ with a functional equation relating $s$ and $1-s$, \emph{RH for $L$} means that
every zero in $0<\Re s<1$ has real part $1/2$. All $L$-functions below are attached to a fixed finite-order
Hecke character, a fixed Hecke character of $\Q(i)$, or a fixed cuspidal automorphic representation.
For these we use the analytic continuation, the functional equation, the order-one Hadamard product of the
completed function, nonvanishing on $\Re s=1$, and the zero count $N(T+1)-N(T)=O(\log(2+T))$. These are
estimates for a fixed function; no uniformity in the field or in the conductor is used. An induced
representation $\Ind\psi$ of a one-dimensional character $\psi$ has $L(s,\Ind\psi)=L(s,\psi)$, a Hecke
$L$-function; no case of Artin's holomorphy conjecture beyond this monomial case is used.

The passage from a one-sided bound to the absence of zeros is always the same. It is proved in
\cref{sec:ana}, and we recall it in the normalization used below.

\begin{lemma}[Landau step]\label{hec:lem:landau}
Let $g\colon[1,\infty)\to\R$ be locally integrable with $g(x)=O(x^{A})$ for some $A$, and let
\[
 G(s)=\int_1^\infty g(x)x^{-s-1}\,\dd x\qquad(\Re s>A).
\]
Let $a>0$. Suppose that $G$ continues meromorphically to $\Re s>a$, that its poles on the real interval
$(a,\infty)$ are simple, and that the residue of $G$ at each of them is positive. If
$g(x)\le Cx^{a}$ for all sufficiently large $x$, then the integral converges for $\Re s>a$, and $G$ is
holomorphic in that half-plane.
\end{lemma}

This is \cref{ana:thm:onesided}, with the pointwise hypothesis of \cref{ana:cor:forms}(a), applied to
$R(x)=g(x)/x$ and $\beta=a$. For a remainder of a logarithmic derivative in the ramp normalization of
\cref{sec:ana} the residues are computed once in \cref{ana:thm:zerotest}.

With $g(x)=xR(x)$ the lemma applies to transforms $\int_1^\infty R(x)x^{-s}\,\dd x$ and bounds
$R(x)\le Cx^{a-1}$. Applied to $-g$, it converts a lower bound into holomorphy when the real residues are
negative. In every application the conclusion is used in the same way: a zero $\rho$ of the tested
$L$-function with $\Re\rho>a$ would produce a pole of $G$ whose residue is a nonzero combination of zero
orders, so there is no such zero.

For the converse implications we use the explicit formula for a fixed $L$-function in the following form.
Under RH for the entire functions $L(s,\chi_1),\dots,L(s,\chi_r)$ of fixed characters, each sum
$\sum_{\Nm\mathfrak p^k\le x}\chi_i(\mathfrak p)^k\log\Nm\mathfrak p$ is $O(\sqrt x\log^2(2x))$, and the
sum for the trivial character is $x+O(\sqrt x\log^2(2x))$ under RH for the corresponding Dedekind zeta
function; see \cite{GrenieMolteni2016Explicit,GrenieMolteni2017Effective} and, for Artin $L$-functions,
\cite[Ch.~3]{Ng2000Limiting}. Prime powers of exponent at least two, inert primes and ramified primes
contribute $O(\sqrt x\log^2(2x))$ to such sums.

\subsection{The cubic field of discriminant \texorpdfstring{$-23$}{-23}}\label{hec:ss:cubic}
{\emergencystretch=\textwidth\sources{grh_transfer_analysis; grh_transfer_analysis_2026-11-25; grh_lean_transfer_analysis; grh_transfers; grh-progression-densities-and-cubic-prime-race; grh_beyond_dirichlet_2026-11-25; grh_method_transfers_2026-11-25_2; grh_dihedral_47_71; grh_transfers_lubin_tate_and_graph_towers}}

Let
\[
 F=\Q(\theta),\qquad \theta^3-\theta-1=0,\qquad K=\Q(\sqrt{-23}),
\]
and let $H$ be the splitting field of $X^3-X-1$. The degree-two Artin $L$-function of this field is a small
nonabelian example with known holomorphy, rather than an arbitrary automorphic family with further
unresolved analytic prerequisites, and it serves as the test case for the transfers discussed in this
article.

\begin{proposition}[The field and its factorizations; \stU]\label{hec:prop:cubicfield}
The group $\Gal(H/\Q)$ is $S_3$ with quadratic subfield $K$, the class group $\Cl(K)$ is cyclic of order
$3$, and $H$ is the Hilbert class field of $K$. If $\psi$ is either nontrivial character of $\Cl(K)$ and
$\rho$ is the two-dimensional irreducible representation of $S_3$, then $L_K(s,\psi)=L(s,\rho)$ and
\begin{align}
\zeta_F(s)&=\zeta(s)L_K(s,\psi),\label{hec:eq:cubic}\\
\zeta_H(s)&=\zeta_K(s)L_K(s,\psi)^2=\zeta(s)L(s,\chi_{-23})L_K(s,\psi)^2.\label{hec:eq:galois}
\end{align}
\end{proposition}

\begin{proof}
The only possible rational roots of $X^3-X-1$ are $\pm1$, and neither is a root. Its discriminant is
$-4(-1)^3-27(-1)^2=-23$. It is squarefree, so the order $\Z[\theta]$ is maximal; it is not a square, so the
transitive Galois group is $S_3$, with quadratic subfield $K$.

Reduced positive definite primitive forms of discriminant $-23$ have first coefficient at most
$\sqrt{23/3}<3$. Enumerating the possibilities gives exactly
\[
(1,1,6),\qquad (2,1,3),\qquad (2,-1,3),
\]
so $h_K=3$ by the correspondence between ideal classes and form classes.

The extension $H/K$ is cyclic of degree $3$. No prime other than $23$ ramifies. Modulo $23$,
\[
X^3-X-1=(X-10)^2(X-3).
\]
The simple root lifts over $\Q_{23}$, so the local decomposition group fixes a root and is contained in a
stabilizer of order two. The quadratic field is ramified at $23$, so inertia is exactly that subgroup of
order two, and its intersection with $A_3=\Gal(H/K)$ is trivial. Thus $H/K$ is unramified at the finite
places, and $K$ has no real places. The degree of $H/K$ equals the class number, so $H$ is the Hilbert
class field.

The permutation representation of $S_3$ on the three roots is $\triv\oplus\rho$, hence
$\zeta_F=\zeta L(s,\rho)$. Inducing either nontrivial character of $A_3$ gives $\rho$: the induced
character takes the values $2,0,-1$ on the identity, the transpositions and the three-cycles.
Compatibility of Artin $L$-functions with induction identifies $L(s,\rho)$ with $L_K(s,\psi)$; the two
nontrivial characters induce isomorphic representations and give the same $L$-function (complex conjugation
of ideals preserves norms and inverts classes). The regular representation of $S_3$ is
$\triv\oplus\hecsgn\oplus\rho^{\oplus2}$, which proves \eqref{hec:eq:galois}.
\end{proof}

The function $L(s)=L_K(s,\psi)$ is the $L$-function of the weight-one cusp form
\begin{equation}\label{hec:eq:eta23}
 f(z)=\eta(z)\eta(23z)=q\prod_{n\ge1}(1-q^n)(1-q^{23n}),\qquad q=e^{2\pi iz},
\end{equation}
of level $23$ and character $\chi_{-23}$
\cite{Zagier2008Elliptic,Serre2003Jordan,KhareLaRosaWiese2023Splitting,LMFDBDatabase}. Its coefficient at
a prime $p\ne23$ is
\[
 a(p)=\tr\rho(\Frob_p)=\#\set{t\in\F_p: t^3-t-1=0}-1\in\set{2,0,-1}.
\]
Put $t=p^{-s}$ and write $D_F(t)$ and $D_\psi(t)$ for the reciprocal Euler factors of $\zeta_F$ and of
$L$ at $p$. All cases, including the ramified prime, are as follows.
\begin{center}
\begin{tabular}{llll}
\toprule
Splitting type of the cubic & $D_F(t)$ & $D_\psi(t)$ & $a(p)$\\
\midrule
$1+1+1$ & $(1-t)^3$ & $(1-t)^2$ & $2$\\
$1+2$ & $(1-t)(1-t^2)$ & $1-t^2$ & $0$\\
$3$ & $1-t^3$ & $1+t+t^2$ & $-1$\\
$p=23$ & $(1-t)^2$ & $1-t$ & $1$\\
\bottomrule
\end{tabular}
\end{center}
At the ramified prime the prime ideal of $K$ has class of order dividing $2$, hence is principal, and its
$\psi$-value is $1$. The completed function
\begin{equation}\label{hec:eq:Lambda23}
 \Lambda(s)=\Bigl(\frac{\sqrt{23}}{2\pi}\Bigr)^{s}\Gamma(s)L(s)
 =23^{s/2}\int_0^\infty f(iy)\,y^{s}\,\frac{\dd y}{y}
\end{equation}
is entire of order one and satisfies $\Lambda(s)=\Lambda(1-s)$. The eta product is positive on the
imaginary axis, so $\Lambda(\sigma)>0$ for every real $\sigma$; in particular $L$ has no real zero in
$(0,1)$, and $L(1/2)\ne0$.

Inside the critical strip all factors in \eqref{hec:eq:cubic} and \eqref{hec:eq:galois} are holomorphic, so
zeros cannot be cancelled by poles. With \cref{src:thm:dgrh} applied to $\zeta$ and $L(s,\chi_{-23})$, the
three assertions RH for $L$, RH for $\zeta_F$ and RH for $\zeta_H$ are therefore equivalent. None of them
is covered by \cref{src:thm:dgrh}.

\paragraph{A positive selector: the principal-class criterion.}
For a nonzero integral ideal $\mathfrak a$ of $K$ put $\Lambda_K(\mathfrak a)=\log\Nm\mathfrak p$ if
$\mathfrak a=\mathfrak p^r$ with $\mathfrak p$ prime and $r\ge1$, and $\Lambda_K(\mathfrak a)=0$ otherwise.
Character orthogonality on the class group gives
$1+\psi(\mathfrak a)+\overline{\psi(\mathfrak a)}=3\,\hecind_{[\mathfrak a]=1}$, and the logarithmic
derivative of \eqref{hec:eq:galois} is, for $\Re s>1$,
\begin{equation}\label{hec:eq:classlogder}
-\frac{\zeta_H'}{\zeta_H}(s)
=3\sum_{[\mathfrak a]=1}\frac{\Lambda_K(\mathfrak a)}{(\Nm\mathfrak a)^s}.
\end{equation}
All coefficients are nonnegative. The sum is over principal prime powers, not over powers of principal
primes: a nonprincipal prime ideal $\mathfrak p$ contributes through $\mathfrak p^3$. The identity can also
be seen locally. If the class of $\mathfrak p$ has order $e\in\set{1,3}$, then $\mathfrak p$ splits into
$3/e$ primes of norm $(\Nm\mathfrak p)^e$ in $H$, so $-\zeta_H'/\zeta_H$ has coefficient
$3\log\Nm\mathfrak p$ at exactly the powers $\mathfrak p^j$ with $e\mid j$; since $H/K$ is unramified, no
Euler factor is exceptional. The same argument, with $1/h$ in place of $1/3$, applies to any unramified
abelian extension of degree $h$. For a rational prime $p$, the coefficient of $p^{-ks}$ in
\eqref{hec:eq:classlogder}, divided by $3\log p$, is
\begin{center}
\begin{tabular}{lll}
\toprule
Type of $p$ & Exponents $k$ with nonzero coefficient & Coefficient$/(3\log p)$\\
\midrule
$\Frob_p=1$ & all $k$ & $2$\\
$\Frob_p$ a transposition & $2\mid k$ & $2$\\
$\Frob_p$ a three-cycle & $3\mid k$ & $2$\\
$p=23$ & all $k$ & $1$\\
\bottomrule
\end{tabular}
\end{center}
The transposition line comes from a prime ideal of norm $p^2$, which is principal because it is generated
by $p$. Define the Euler--Kronecker constant
$\gamma_H=\lim_{s\to1}\bigl(\zeta_H'(s)/\zeta_H(s)+1/(s-1)\bigr)$ and the normalized smoothed remainder
\begin{equation}\label{hec:eq:RH}
\mathcal R_H(x)=
3\!\!\sum_{\substack{\Nm\mathfrak a\le x\\{}[\mathfrak a]=1}}
\Lambda_K(\mathfrak a)
\Bigl(\frac1{\Nm\mathfrak a}-\frac1x\Bigr)
-\log x+(1+\gamma_H).
\end{equation}
The sign and the constant are forced by the Laurent expansion at $s=1$; another additive constant would
destroy a decaying bound.

\begin{theorem}[Principal-class criterion; \stR]\label{hec:thm:principal}
Fix $0<\beta<1$. If there is a constant $C$ with
\begin{equation}\label{hec:eq:principalbound}
\mathcal R_H(x)\le Cx^{\beta-1}\qquad\text{for all sufficiently large }x,
\end{equation}
then $\zeta_H$ has no zero with $\Re s>\beta$, and consequently neither do $\zeta_K$, $\zeta_F$ and
$L_K(s,\psi)$. If \eqref{hec:eq:principalbound} holds for every $\beta>1/2$, then RH holds for $\zeta_H$,
$\zeta_F$ and $L_K(s,\psi)$.
\end{theorem}

\begin{proof}
By \eqref{hec:eq:classlogder}, $\mathcal R_H$ is the remainder \eqref{ana:eq:RK} of the Dedekind zeta
function of $H$: the sum is $\sum_{n\le x}\Lambda_H(n)(1/n-1/x)$, where $\Lambda_H(n)$ groups the
coefficients of $-\zeta_H'/\zeta_H$ by rational norm, and the constant is $1+\gamma_H$. With $u=s-1$ one has
$-\zeta_H'/\zeta_H(1+u)=1/u-\gamma_H+O(u)$, and division by $s(s-1)=u(1+u)$ gives
$u^{-2}-(1+\gamma_H)u^{-1}+O(1)$; this is the normalization that makes the Mellin transform of
$\mathcal R_H$ regular at $s=1$ (\cref{ana:lem:mellin}). At a zero $\rho$ of $\zeta_H$ of multiplicity
$m_\rho$ the transform has a simple pole with residue $m_\rho/(\rho(1-\rho))$, even if the zero is
multiple. \Cref{ana:thm:zerotest}, in the instance \cref{ana:ex:instances}(a) with $K=H$, therefore excludes
every zero of $\zeta_H$ in $\Re s>\beta$. The factorizations \eqref{hec:eq:cubic} and \eqref{hec:eq:galois},
with all factors holomorphic in the strip, give the statements for the other functions, and the
functional equations give the zeros to the left of the critical line.
\end{proof}

The theorem makes no use of \cref{src:thm:dgrh}: the estimate would give the Riemann hypothesis for every
factor of $\zeta_H$ at once. Because the selector is the full class-character product, every residue has the same
sign, and no auxiliary zero-avoiding twist is required. By the averaged forms of the Landau step in
\cref{ana:cor:forms,ana:thm:moving}, the pointwise bound may be replaced, for every $\eps>0$, by any one of
\[
 \int_X^{2X}(\mathcal R_H)_+\,\dd x\ll_\eps X^{1/2+\eps},\qquad
 \int_X^{2X}(\mathcal R_H)_+^2\,\dd x\ll_\eps X^{\eps},\qquad
 \int_X^{2X}\mathcal R_H\,\dd x\le C_\eps X^{1/2+\eps},
\]
the first two on dyadic $X$ and the third for all large real $X$; since
$\mathcal R_H(x)=O(\log^2(2x))$, integer starting points suffice in the third. Conversely, RH for this
fixed field supplies the conventional smoothed error bounds, and hence these averaged bounds. The three
targets are therefore weaker forms of control to aim for, not conjectures logically weaker than RH for
$\zeta_H$. The local normalizations
were checked by exact computation: the coefficients of \eqref{hec:eq:eta23} agree with the Euler product
obtained from the splitting of the cubic through $n=2000$ (303 primes), and \eqref{hec:eq:classlogder} was
checked on twelve powers of each of these primes
(\artifact{grh_beyond_dirichlet_2026-11-25/code/conductor23_benchmark.py},
\artifact{grh_lean_transfer_analysis/computation/cubic23_exact.py}).
{\emergencystretch=14em\par}

\paragraph{A twisted coefficient system.}
The source proof works with coefficients weighted by $1-\chi_d$ for an auxiliary quadratic character.
The same shape is available here, at the price of a second hypothesis. Write $\epsilon$ for the sign
character of $S_3$. The character table gives $\abs{\tr\rho(g)}^2=1+\epsilon(g)+\tr\rho(g)$ on every class,
so
\[
 \rho\otimes\rho^*\simeq\triv\oplus\epsilon\oplus\rho ,
\]
and $\abs{\tr\rho(g^k)}^2\ge0$ is the character of this representation at every prime-power argument
(a transposition gives $0$ for odd $k$ and $4$ for even $k$; a three-cycle gives $1$ unless $3\mid k$, when
it gives $4$). Let $d\equiv7\pmod8$ be a prime, $d\ne23$, let $\chi_d(n)=\legendre{n}{d}$, and put
$S_d=\set{23,d}$. With a superscript $S_d$ denoting omission of the Euler factors at those primes, define
\begin{align*}
 P_d(s)&=\zeta^{S_d}(s)L^{S_d}(s,\epsilon)L^{S_d}(s,\rho),\\
 Q_d(s)&=L^{S_d}(s,\chi_d)L^{S_d}(s,\epsilon\chi_d)L^{S_d}(s,\rho\otimes\chi_d).
\end{align*}
Neither one-dimensional character in $Q_d$ is principal, and $\rho\otimes\chi_d$ is induced from the
character $\psi\cdot\chi_d|_{G_K}$, whose cubic part cannot be cancelled by a quadratic character. Thus
$Q_d$ is entire and nonzero at $s=1$, while $P_d$ has a simple pole there. For $n=p^k$ with $p\notin S_d$
put
\begin{equation}\label{hec:eq:lambdad}
 \lambda_d(p^k)=(\log p)\,\abs{\tr\rho(\Frob_p^k)}^2\bigl(1-\chi_d(p)^k\bigr),
\end{equation}
and $\lambda_d(n)=0$ otherwise. Then $0\le\lambda_d(n)\le8\Lambda(n)$ and
$\sum_n\lambda_d(n)n^{-s}=-P_d'/P_d(s)+Q_d'/Q_d(s)$ for $\Re s>1$. The pole coefficient at $1$ is one,
the mean of the squared trace. Let $A_d$ be the entire function $(s-1)P_d(s)$, put
$G_d(s)=Q_d'(s)/Q_d(s)-A_d'(s)/A_d(s)-1$ and $c_d=G_d(1)\in\R$, and define
\begin{equation}\label{hec:eq:Rd}
 R_d(x)=\sum_{n\le x}\lambda_d(n)\Bigl(\frac1n-\frac1x\Bigr)-\log x-c_d.
\end{equation}

\begin{proposition}[Twisted target; \stR]\label{hec:prop:twisted}
Fix $1/2<\beta<1$. Suppose that for every $T\ge0$ there is a prime $d\equiv7\pmod 8$, $d\ne23$, such that
\begin{enumerate}[label=\textup{(\Alph*)}]
\item $R_d(x)\le C_dx^{\beta-1}$ for all sufficiently large $x$, for some finite $C_d$, and
\item $Q_d(s)\ne0$ whenever $\Re s\ge\beta$ and $\abs{\Im s}\le T$.
\end{enumerate}
Then $\zeta(s)L(s,\epsilon)L(s,\rho)$, and in particular $L(s,\rho)$, has no zero in $\Re s>\beta$.
\end{proposition}

\begin{proof}
The sum of the series is $s/(s-1)+G_d(s)$, so the same integration as before gives
\[
 \Phi_d(s)=\int_1^\infty R_d(x)x^{-s}\,\dd x=\frac{G_d(s)/s-c_d}{s-1}.
\]
The point $s=1$ is removable, and at $\omega\ne0,1$
\begin{equation}\label{hec:eq:resd}
 \Res_{s=\omega}\Phi_d(s)=\frac{\ord_\omega A_d-\ord_\omega Q_d}{\omega(1-\omega)}.
\end{equation}
All the representations have finite image, so the removed Euler polynomials vanish only on $\Re s=0$;
omitting $S_d$ changes no zero order in $\Re s>0$ (it does change the constant $c_d$). The bound on
$\lambda_d$ gives $R_d(x)=O_d(1+\log^2x)$. Suppose $\omega$ is a zero of $\zeta L(s,\epsilon)L(s,\rho)$ with
$\Re\omega>\beta$, take $T\ge\abs{\Im\omega}$, and use the same $d$ in (A) and (B). For real $\beta<t<1$
the residue \eqref{hec:eq:resd} is $\ord_tA_d/(t(1-t))\ge0$ by (B), and there is no real pole on
$[1,\infty)$. By \cref{hec:lem:landau}, applied to $g(x)=xR_d(x)$, the function $\Phi_d$ is holomorphic in
$\Re s>\beta$. But at $\omega$ the target order is positive and the auxiliary order vanishes by (B), so
\eqref{hec:eq:resd} is nonzero, a contradiction.
\end{proof}

If (A) and (B) hold for every $\beta>1/2$, the functional equation gives RH for $L(s,\rho)$. With
\cref{src:thm:dgrh} the one-dimensional factors of $Q_d$ are already zero-free in $\Re s>1/2$; the new
separation problem is the degree-two factor $L(s,\rho\otimes\chi_d)$, and (A) and (B) must hold for a
common $d$. Unlike \cref{hec:thm:principal}, this route needs an auxiliary nonvanishing statement, because
the residue \eqref{hec:eq:resd} is a difference of zero orders.

\paragraph{A sign test on prime root counts.}
The strongest form of the reduction uses primes only, a single cutoff, a sign rather than a rate, and a
sparse set of cutoffs. Consider the first Riesz mean of the trace over primes,
\begin{equation}\label{hec:eq:T}
 \mathcal{T}(x)=\sum_{p\le x,\ p\ne23}a(p)\log p\,\Bigl(1-\frac px\Bigr).
\end{equation}
Call an increasing unbounded sequence $(x_j)$ \emph{admissible} if
$x_{j+1}-x_j=O(x_j^{1-\delta})$ for some fixed $0<\delta\le1$; for instance $x_j=j^k$ with $\delta=1/k$,
for any integer $k\ge1$.

The completed function \eqref{hec:eq:Lambda23} is the one used in \cref{zer:sec:window}, and the zero sum
that controls the oscillation is
\[
 Z=\sum_{\rho}\frac1{1/4+(\Im\rho)^2},
\]
over the nontrivial zeros of $L$ with multiplicity, as in \eqref{zer:eq:Z}. The curvature bound for positive
eta products, \cref{zer:prop:eta}, gives in the case $\eta(z)\eta(23z)$ the following
(\cref{zer:cor:cubicZ}): if RH holds for $L$, then
$Z=2\Lambda'(1)/\Lambda(1)\le6/(\pi\sqrt{23})$.

\begin{theorem}[Sparse sign test; \stR]\label{hec:thm:sign}
Let $(x_j)$ be an admissible sequence. Then RH for $L$ is equivalent to
\[
 \mathcal{T}(x_j)\le0\qquad\text{for all sufficiently large }j.
\]
Moreover, RH for $L$ implies
\begin{equation}\label{hec:eq:limsup}
 \limsup_{x\to\infty}\frac{\mathcal{T}(x)}{\sqrt x}\le\frac6{\pi\sqrt{23}}-\frac23<-0.268 .
\end{equation}
\end{theorem}

\begin{theorem}[Oscillation on sparse sequences; \stU]\label{hec:thm:osc}
Let $(x_j)$ be an admissible sequence.
\begin{enumerate}[label=\textup{(\roman*)}]
\item $\mathcal{T}(x_j)<0$ for infinitely many $j$, and
$\liminf_{j}\mathcal{T}(x_j)/\sqrt{x_j}\le6/(\pi\sqrt{23})-2/3$.
\item If $L$ has a zero with real part $\beta>1/2$, then for every $1/2\le a<\beta$,
\[
 \limsup_{j\to\infty}\frac{\mathcal{T}(x_j)}{x_j^{a}}=+\infty,\qquad
 \liminf_{j\to\infty}\frac{\mathcal{T}(x_j)}{x_j^{a}}=-\infty .
\]
\item In all cases
\[
 \liminf_{j\to\infty}\frac{\mathcal{T}(x_j)}{\sqrt{x_j}}<-\frac23<\limsup_{j\to\infty}\frac{\mathcal{T}(x_j)}{\sqrt{x_j}}.
\]
\end{enumerate}
\end{theorem}

\begin{proof}[Proof of \cref{hec:thm:sign,hec:thm:osc}]
\emph{The explicit formula and the bias from squares.}
Write $b(p^j)=\tr(\rho(\Frob_p)^j)\log p$ at unramified primes, with the inertia-invariant trace at $23$.
Assume RH for $L$. Mellin inversion on $\Re s=2$ with the kernel $x^s/(s(s+1))$, which inverts to
$(1-p^j/x)_+$, expresses the Riesz mean of $b$ as the integral of $(-L'/L)(s)x^s/(s(s+1))$. Shift the line
to $\Re s=-1/2$ along heights $T$ between successive integers at distance $\gg1/\log T$ from the zero
ordinates, where the logarithmic derivative is $O(\log^2(2+\abs s))$; such heights exist by the local zero
count, and the two horizontal integrals are $O_x(T^{-2}\log^2T)$. On the new line the functional equation
gives $O(\log(2+\abs s))$ and the integral is $O(x^{-1/2})$. At $s=0$ the pole of the gamma factor and
$\Lambda(0)>0$ give a simple zero of $L$, so the double pole there contributes $O(1+\log x)$. Hence
\begin{equation}\label{hec:eq:riesz}
 \sum_{p^j\le x}b(p^j)\Bigl(1-\frac{p^j}x\Bigr)
 =-\sum_{\rho}\frac{x^{\rho}}{\rho(\rho+1)}+O\bigl((1+\log x)^2\bigr),
\end{equation}
the zero sum being absolutely convergent. Since $\abs{\rho+1}\ge\abs\rho$ on the critical line, the zero
sum has modulus at most $Z\sqrt x$. At the squares, the character $h\mapsto\tr\rho(h^2)$ has group mean
$(2+3\cdot2+2\cdot(-1))/6=1$, so the Chebotarev theorem for the fixed field $H$ and partial summation give
\begin{equation}\label{hec:eq:squares}
 \sum_{p\le\sqrt x,\ p\ne23}b(p^2)\Bigl(1-\frac{p^2}x\Bigr)
 =\Bigl(\int_0^1(1-y^2)\,\dd y+o(1)\Bigr)\sqrt x=\Bigl(\frac23+o(1)\Bigr)\sqrt x
\end{equation}
without any hypothesis on zeros. The terms of exponent at least three contribute $O(x^{1/3}\log^2x)$, and
the powers of $23$ contribute $O(\log x)$. Subtracting from \eqref{hec:eq:riesz} and using the bound
$Z\le6/(\pi\sqrt{23})$ of \cref{zer:cor:cubicZ} proves \eqref{hec:eq:limsup}, and with it the implication
from RH to the eventual sign, for every sequence tending to infinity.

\emph{From a one-sided bound to zeros.}
For $\Re s>1$,
\begin{equation}\label{hec:eq:Tmellin}
 \int_1^\infty \mathcal{T}(x)x^{-s-1}\,\dd x=\frac{-L'(s)/L(s)-E(s)}{s(s+1)},
\end{equation}
where $E$ collects the powers of exponent at least two and the ramified prime, and is holomorphic in
$\Re s>1/2$. A zero $\rho$ of $L$ of multiplicity $m$ with $\Re\rho>1/2$ gives a simple pole with residue
$-m/(\rho(\rho+1))\ne0$, and by positivity of $\Lambda$ on the real axis there is no real pole in
$(1/2,\infty)$. By \cref{ana:thm:onesided} with \cref{ana:cor:forms}(a), applied to $\mathcal{T}(x)/x$ and to
$-\mathcal{T}(x)/x$ with $\beta=a$, a bound $\mathcal{T}(x)\le Cx^{a}$ or $-\mathcal{T}(x)\le Cx^{a}$ with $a>1/2$, valid for all
large $x$, excludes all zeros in $\Re s>a$.

\emph{Interpolation between the sample points.}
Suppose $L$ has no zero in $\Re s>\theta$, where $1/2\le\theta\le1$. Then
$U(x)=\sum_{p\le x,\,p\ne23}a(p)\log p$ is $O_\eps(x^{\theta+\eps})$ for each $\eps>0$: at $\theta=1$ this
is trivial, and otherwise it follows from the truncated explicit formula for the fixed Artin factor at
half-integers with a height $T\asymp x^2$ chosen away from the ordinates, in which the zeros contribute
$O(x^\theta\sum_{\abs{\Im\rho}\le T}1/\abs\rho)=O(x^\theta\log^2T)$, the truncation error is
$O(\log^{O(1)}x)$, and the omitted prime powers cost $O(\sqrt x\log^2x)$
\cite[\S3.4 and Example~3.5.6]{Ng2000Limiting}. Between primes,
\[
 \abs{\mathcal{T}'(x)}=x^{-2}\Bigl\lvert\sum_{p\le x,\ p\ne23}p\,a(p)\log p\Bigr\rvert=O_\eps(x^{\theta-1+\eps})
\]
by partial summation, and $T$ is continuous at primes. If $\mathcal{T}(x_j)\le Cx_j^{a}$ for all large $j$, with
$a\ge1/2$, then integrating across a gap of length $O(x_j^{1-\delta})$ gives the upper bound
$\mathcal{T}(x)\ll_\eps x^{a}+x^{\theta-\delta+\eps}$ for all large $x$, and the previous step excludes zeros in
$\Re s>\max(a,\theta-\delta)$. Starting from $\theta=1$ and iterating at most $\lceil1/(2\delta)\rceil$
times shows that $L$ is zero-free in $\Re s>a$. The same holds for $-\mathcal{T}$.

\emph{Conclusions.}
With $C=0$ and $a=1/2$ the iteration shows that $\mathcal{T}(x_j)\le0$ for all large $j$ implies that $L$ has no
zero in $\Re s>1/2$, and the functional equation gives RH for $L$. This completes the proof of
\cref{hec:thm:sign}. If $\mathcal{T}(x_j)$ were nonnegative for all large $j$, the same argument for $-\mathcal{T}$ would give
RH for $L$, hence \eqref{hec:eq:limsup} and eventual strict negativity, a contradiction; so $\mathcal{T}(x_j)<0$
infinitely often. If $L$ has a zero with real part $\beta>1/2$ and $1/2\le a<\beta$, a finite upper bound
for $\mathcal{T}(x_j)/x_j^{a}$ or for $-\mathcal{T}(x_j)/x_j^{a}$ would make $L$ zero-free in $\Re s>a$; this proves (ii).
Combining (ii) at $a=1/2$ with \eqref{hec:eq:limsup} gives the second assertion of (i).

For (iii) we may assume RH for $L$, since otherwise (ii) gives both infinite extrema. With $t=\log x$,
\eqref{hec:eq:riesz} and \eqref{hec:eq:squares} give
\[
 e^{-t/2}\mathcal{T}(e^t)=-\frac23+\Phi(t)+o(1),\qquad
 \Phi(t)=-\sum_{\rho=1/2+i\gamma}\frac{e^{i\gamma t}}{\rho(\rho+1)} .
\]
Absolute uniform convergence makes $\Phi$ real and uniformly continuous. Positivity of $\Lambda(1/2)$ rules
out $\gamma=0$, so the average of $\Phi$ over $[0,T]$ tends to zero. There is at least one zero: a
zero-free entire function of order one is the exponential of an affine function, the symmetry would make
it constant, and this contradicts Stirling's formula at large positive real $s$. Choose one positive
ordinate $\gamma$, of multiplicity $m$, and put $C=-m/(\rho(\rho+1))$, the nonzero Fourier coefficient of
$\Phi$ at $\gamma$. The weight $1+\cos(\gamma t+\phi)$ is nonnegative with mean $1$, and its product with
$\Phi$ has mean $\Re(e^{-i\phi}C)$. The phases giving $-\abs C$ and $\abs C$ show
$\liminf\Phi\le-\abs C$ and $\limsup\Phi\ge\abs C$. Finally $\log x_{j+1}-\log x_j\to0$, and uniform
continuity gives the same two extrema along $(x_j)$.
\end{proof}

The script \artifact{grh-progression-densities-and-cubic-prime-race/code/check_eta23.py} compares the
78{,}497 unramified primes below $10^6$ with the coefficients of the eta product by two independent exact
algorithms, and gives $\mathcal{T}(10^6)/1000=-0.788953233$ to the displayed precision. No finite
computation bears on the eventual sign.
{\emergencystretch=14em\par}

\begin{remark}[Two-scale form]\label{hec:rem:twoscale}
Let $B(y)=\sum_{p\le y,\,p\ne23}a(p)\log p$ and, for a dilation $h>1$,
$J_h(x)=-\int_x^{hx}B(y)\,y^{-2}\,\dd y$, which is the statistic \eqref{zer:eq:Jh} for the weight $a(p)$.
\Cref{zer:thm:cubicrace} shows that RH for $L$ is equivalent to $J_4(x)\ge0$ for all sufficiently large $x$,
and to $\liminf_{x\to\infty}\sqrt x\,J_4(x)\ge1-9/(\pi\sqrt{23})>0.4026$, and that $J_4$ takes positive
values arbitrarily far out unconditionally. \Cref{hec:thm:sign} replaces the two continuous cutoffs by a
single one on a prescribed sparse sequence.
\end{remark}

\paragraph{Irreducible primes.}
The geometric part of the source proof produces nonnegative local amounts. A nonnegative class function
on a group $G$ whose coefficients on all nonlinear characters are negative is positive on every nonidentity
element of $[G,G]$, and the indicator of these elements is the pointwise minimal such function in the
normalization of \cref{sec:nab}. For $S_3$ this indicator is that of the three-cycles, which asks whether
the cubic is irreducible modulo $p$.

\begin{corollary}[Irreducible-prime criterion for $S_3$ cubics; \stR]\label{hec:thm:irred}
Let $P\in\Z[X]$ be a monic irreducible cubic with nonsquare discriminant, let $E$ be its splitting field,
with group $S_3$, and let $\tau$ be the standard representation. Assume \cref{src:thm:dgrh}. Then RH for
$L(s,\tau)$ is equivalent to RH for $\zeta_E$, and hence to each of the equivalent statements of
\cref{nab:thm:nonsplit} with $N=[S_3,S_3]=A_3$. For this choice the set $\mathcal P$ consists of the primes
$p\nmid\disc(P)$ at which $P\bmod p$ is irreducible, $\delta=1/2-1/6=1/3$, and the weight is
\[
 3\,\hecind_{A_3\setminus\set1}=1+\hecsgn-\chi_\tau .
\]
In particular RH for $L(s,\tau)$ is equivalent to the bound, for every $\eps>0$,
\begin{equation}\label{hec:eq:irredP}
 \sum_{\substack{p\le x,\ p\nmid\disc(P)\\ P\bmod p\ \mathrm{irreducible}}}\log p-\frac x3\le C_{P,\eps}x^{1/2+\eps}.
\end{equation}
\end{corollary}

\begin{proof}
In the order identity, transposition, three-cycle, the characters are $\chi_\tau=(2,0,-1)$ and
$\hecsgn=(1,-1,1)$, which gives the displayed weight; a prime $p\nmid\disc(P)$ has a three-cycle as
Frobenius exactly when $P\bmod p$ is irreducible. The representation $\tau$ is induced from a nontrivial
character of $A_3$, so $L(s,\tau)$ is an entire Hecke $L$-function, and the regular representation gives
$\zeta_E=\zeta\,L(s,\hecsgn)L(s,\tau)^2$ with all factors holomorphic in the critical strip. With
\cref{src:thm:dgrh} for the first two factors, RH for $\zeta_E$ and RH for $L(s,\tau)$ are equivalent, and
\eqref{hec:eq:irredP} is statement (iv) of \cref{nab:thm:nonsplit}.
\end{proof}

The first Riesz mean of the same primes gives a two-sided form without loss of $\eps$.

\begin{proposition}[Riesz form of the irreducible-prime criterion; \stR]\label{hec:prop:irredriesz}
With $P$ and $\tau$ as in \cref{hec:thm:irred}, put
\[
 I_P(x)=3\sum_{\substack{p\le x,\ p\nmid\disc(P)\\ P\bmod p\ \mathrm{irreducible}}}
 \log p\,\Bigl(1-\frac px\Bigr).
\]
Assume \cref{src:thm:dgrh}. The following are equivalent.
\begin{enumerate}[label=\textup{(\arabic*)}]
\item All nontrivial zeros of $L(s,\tau)$ have real part $1/2$.
\item $I_P(x)=x/2+O_P(\sqrt x)$.
\item For each $\eps>0$ and all sufficiently large $x$, $I_P(x)\le x/2+C_{P,\eps}x^{1/2+\eps}$.
\end{enumerate}
\end{proposition}

\begin{proof}
Let $w_*$ be the indicator of the three-cycles and $U(x)$ the sum over prime powers $p^k\le x$,
$p\nmid\disc(P)$, of $3w_*(g_p^k)\log p\,(1-p^k/x)$, where $g_p$ is a Frobenius element. Chebyshev's bound
gives $U-I_P=O(\sqrt x)$. For $\Re s>1$,
\[
 \int_1^\infty U(x)x^{-s-1}\,\dd x=\frac{A(s)}{s(s+1)},\qquad
 A(s)=-\frac{\zeta'}{\zeta}(s)-\frac{L'}L(s,\hecsgn)+\frac{L'}L(s,\tau)+B(s),
\]
where $B$ accounts for the finitely many excluded Euler factors and is holomorphic in $\Re s>0$. The
residue at $s=1$ is $1/2$. After subtraction of $1/(2(s-1))$, the transform of $U(x)-x/2$ has, at a zero
$\beta$ of $L(s,\tau)$ with $\Re\beta>1/2$ of multiplicity $m_\beta$, a simple pole with residue
$m_\beta/(\beta(\beta+1))$, positive at real $\beta$; by \cref{src:thm:dgrh} the first two terms of $A$
contribute no other pole in this region. Assume (3) and fix $a=1/2+\eps<1$. By \cref{hec:lem:landau} the
transform is holomorphic in $\Re s>a$, so $L(s,\tau)$ has no zero there; let $\eps\to0$ and use the
functional equation. This proves (1).

Assume (1). All three factors then satisfy their critical-line statements. Start from Mellin inversion
on $\Re s=2$ with the full logarithmic derivatives, adding back the omitted local terms at a cost
$O_P(\log x)$, and shift to $\Re s=-1/2$ along heights at distance $\gg_P1/\log T$ from the ordinates. The
logarithmic derivatives are $O_P(\log^2T)$ on the horizontal segments and $O_P(\log(2+\abs t))$ on the new
line, the residues at the nontrivial zeros have absolute sum $O_P(\sqrt x)$ because
$\sum_\rho(1+\abs{\Im\rho})^{-2}$ converges, and the possible double pole at $0$ costs $O_P(1+\log x)$.
This gives $U(x)=x/2+O_P(\sqrt x)$, which is (2), and (2) implies (3).
\end{proof}

For $P=X^3-X-1$, with $E=H$, the estimates in \cref{hec:thm:irred,hec:prop:irredriesz} are open. With
$I(p)=1$ if $X^3-X-1$ is irreducible modulo $p$ and $I(p)=0$ otherwise, \eqref{hec:eq:irredP} reads
\begin{equation}\label{hec:eq:irred23}
 \sum_{p\le x,\ p\ne23}I(p)\log p-\frac x3\le C_\eps x^{1/2+\eps}\qquad\text{for every }\eps>0 ,
\end{equation}
and under \cref{src:thm:dgrh} it is equivalent to RH for $\zeta_H$.

\paragraph{Obstructions to an automatic Dirichlet reduction.}
The weights $a(p)$ and $I(p)$ cannot be produced by changing the modulus in the source construction, and
the field $H$ cannot be substituted for the auxiliary field.

\begin{proposition}[Finite-modulus obstruction; \stU]\label{hec:prop:modulus}
For every integer $q\ge1$, infinitely many unramified primes $p\equiv1\pmod q$ have $a(p)=2$, and
infinitely many have $a(p)=-1$. Consequently $a(p)$ is not a function of $p$ modulo any fixed integer,
even after deleting finitely many primes, and it is not a finite linear combination of Dirichlet
characters. No finite product or quotient of Dirichlet $L$-functions has the Euler product of $L$.
\end{proposition}

\begin{proof}
Put $C=\Q(\zeta_q)$. The field $H\cap C$ is abelian over $\Q$, and the largest abelian quotient of $S_3$
is its sign quotient, so $H\cap C\subseteq K$. The identity and a three-cycle act trivially on this
intersection, so the pairs $(1,1)$ and $((123),1)$ define elements of $\Gal(HC/\Q)$, viewed as a fibre
product over $\Gal(H\cap C/\Q)$. The Chebotarev theorem gives infinitely many primes for each of their
conjugacy classes. The cyclotomic component is trivial, so $p\equiv1\pmod q$, while the traces are $2$ and
$-1$. A finite list of Dirichlet characters has a common modulus, on whose identity class every character
equals $1$, and no fixed linear combination can take both values there. The last assertion follows by
comparing the prime coefficients of logarithmic derivatives.
\end{proof}

The mean-square form of this obstruction is the case $E=H$ of \cref{nab:thm:periodic}. For a modulus $m$
the field $H\cap\Q(\zeta_m)$ is abelian over $\Q$, hence is $\Q$ or $K$, and it is $K$ exactly when
$23\mid m$, because $K$ has conductor $23$; so the subgroup $H_m$ of \cref{nab:thm:periodic} is $S_3$ if
$23\nmid m$ and $A_3$ if $23\mid m$. The weights expand as $a=\chi_\rho$ and
$I=\tfrac13(1+\epsilon-\chi_\rho)$ in the characters of $S_3$. Reading off \eqref{nab:eq:energy} and
\cref{nab:prop:error} gives the following explicit values.

\begin{corollary}[Mean-square distance to congruence weights; \stU]\label{hec:prop:meansq}
Let $m\ge1$.
\begin{enumerate}[label=\textup{(\roman*)}]
\item For every real function $w$ on $(\Z/m)^\times$,
\[
 \lim_{X\to\infty}\frac1{\pi(X)}\sum_{\substack{p\le X\\p\nmid23m}}\bigl(a(p)-w(p\bmod m)\bigr)^2
 =1+\frac1{\varphi(m)}\sum_{r\in(\Z/m)^\times}w(r)^2 ,
\]
so the best periodic approximation to $a(p)$ is $w=0$: the coset average of the nonlinear character
$\chi_\rho$ vanishes and $\norm{\chi_\rho}^2_{S_3}=1$.
\item The least value over $w$ of
$\lim_{X\to\infty}\pi(X)^{-1}\sum_{p\le X,\,p\nmid23m}(I(p)-w(p\bmod m))^2$ is
$\tfrac19+\tfrac19=\tfrac29$ if $23\nmid m$, where the coefficients of $\epsilon$ and $\chi_\rho$ are lost,
and $\tfrac19$ if $23\mid m$, where only that of $\chi_\rho$ is lost. The best $\set{0,1}$-valued
congruence rule, with the modulus freely chosen, misclassifies a proportion
$\tfrac16\bigl(\min\set{2,1}+\min\set{0,3}\bigr)=\tfrac16$ of the primes: it predicts irreducibility on the
classes in which $p$ splits in $K$.
\end{enumerate}
\end{corollary}

These values exclude an exact replacement of the weight by a fixed collection of Dirichlet characters;
they do not exclude cancellation in the signed sum \eqref{hec:eq:irred23}, nor constructions whose modulus grows
with the cutoff (\cref{nab:rem:fixed}).

\begin{remark}[Structural obstructions]\label{hec:rem:structural}
Three further facts constrain a transfer of the source construction to this field.
First, although $H$ contains the imaginary quadratic field $K$ and is totally imaginary, it is not a CM
field: in a Galois CM field complex conjugation is central, whereas in $S_3$ it is a noncentral
transposition. For the same reason $H$ is contained in no Galois CM field, since a central element would
restrict centrally to $H$; and it is contained in no abelian field. Replacing the auxiliary CM field of the
source by $H$ is therefore not possible, and an extension has to use relative or adelic class information
over the CM base $K$. Second, the new information sits at the split primes: if $p$ is inert in $K$, the
prime ideal above it is $(p)$, which is principal, so $\psi((p))=1$, and the two nontrivial class characters
are distinguished from the trivial one only at split prime ideals, whereas the divisibility construction
of the source is organized around primes at which the auxiliary quadratic character is not $1$. Third, the
one-dimensional representations of $S_3$ are $\triv$ and $\epsilon$, and every virtual combination of them
takes the same value at the identity and at a three-cycle, while $\tr\rho$ takes the values $2$ and $-1$;
the degree-two factor is thus not reached by any argument through one-dimensional representations of
$\Gal(H/\Q)$. These are limitations of particular transfer routes, not statements of logical independence
from \cref{src:thm:dgrh}.
\end{remark}

\paragraph{Controlling a degree-two zero sum.}
For $K=\Q[\alpha]$, $\alpha^3-\alpha-1=0$, the standard Artin factor is the Hecke
$L$-function of the weight-one form
\[
 f(z)=\eta(z)\eta(23z),\qquad\zeta_K(s)=\zeta(s)L(s,\rho),
\]
with completed function $\zerLL(s)=(\sqrt{23}/(2\pi))^s\Gamma(s)L(s,\rho)$ and
$\zerLL(s)=\zerLL(1-s)$ \cite{Serre2003Jordan,Zagier2008Elliptic}. The eta product is
positive on the positive imaginary axis, and its Mellin integral makes $\zerLL(s)>0$ for
all real $s$, including $s=1/2$. The zero sum $Z$ of \eqref{zer:eq:Z} needs quantitative
control of the logarithmic derivative as well.

\begin{proposition}[\stU]\label{zer:prop:eta}
Let $F(y)=\prod_j\eta(id_jy)^{r_j}$ with finitely many $d_j,r_j>0$. For the entire Mellin
transform $M(s)=\int_0^\infty F(y)y^s\,\dd y/y$ let
\[
 c=\frac\pi{12}\inf_{y>0}\sum_jr_j\max\set{d_jy,(d_jy)^{-1}}.
\]
For every real $\sigma$, $(\log M)''(\sigma)\le1/c$. The transform has no zero at any $s$
with $\abs{\Im s}\le(\pi/2)\sqrt c$.
\end{proposition}
\begin{proof}
Set $v=\log y$. The product formula for $\eta(ie^v)$ starts with $\exp(-\pi e^v/12)$. For
$z=ae^v$, the derivative of $\log(1-e^{-z})$ in $v$ is $z/(e^z-1)$, decreasing in $v$. Thus
$-(\log\eta(ie^v))''\ge(\pi/12)e^v$. Termwise differentiation is justified locally by
exponential convergence. Eta inversion gives $\eta(ie^{-v})=e^{v/2}\eta(ie^v)$, making this
second logarithmic derivative even. Apply the two lower bounds to each factor separately.
The potential $V(v)=-\log F(e^v)$ satisfies $V''\ge c$. The eta product and its inversion
also give double exponential decay in $v$ at both ends, up to an exponential linear factor.
The integral is entire and all differentiations below are justified.

Give $U$ density proportional to $\exp(\sigma v-V(v))$, and write $\mu=\mathbb E\,U$.
Logarithmic differentiation gives $(\log M)''(\sigma)=\operatorname{Var}U$. Integration by
parts gives $\mathbb E[(U-\mu)(V'(U)-\sigma)]=1$. Subtracting $V'(\mu)$ in place of $\sigma$
leaves this expectation unchanged. Strong convexity then gives $1\ge c\operatorname{Var}U$.
Finally, $\cos w\ge\pi/4-w^2/\pi$ on $\R$, strictly except at $w=\pm\pi/2$: on $[0,\pi/2]$
this follows by differentiating and using $\sin w\ge2w/\pi$, and past $\pi/2$ one has
$2w/\pi\ge1$. The centered characteristic function consequently has strictly positive real
part whenever $t^2\operatorname{Var}U\le\pi^2/4$. It equals
$e^{-it\mu}M(\sigma+it)/M(\sigma)$.
\end{proof}

\begin{corollary}\label{zer:cor:cubicZ}
For $f=\eta(z)\eta(23z)$ one has $c=\pi\sqrt{23}/6$, and the completed function
$\zerLL(s)$ has no zero with $\abs{\Im s}\le(\pi/2)\sqrt c=2.48915\ldots$ \textup{(\stU)}.
If all zeros of $L(s,\rho)$ in the critical strip have real part $1/2$, then
$Z\le1/c=6/(\pi\sqrt{23})$ and
\[
 \liminf_{x\to\infty}\sqrt x\,J_4(x)\ \ge\ 1-\frac9{\pi\sqrt{23}}>0.4026
\]
for the cubic weight $b(p)=\#\set{y\in\F_p:y^3-y-1=0}-1$.
\end{corollary}
\begin{proof}
For $d_1=1$, $d_2=23$ the infimum is $2\sqrt{23}$, attained at $y=1/\sqrt{23}$, by
comparison with $1/y+23y$. Multiplying $M(s)$ by $23^{s/2}$ gives $\zerLL(s)$ and leaves
the second logarithmic derivative unchanged. The symmetry forces $(\log\zerLL)'(1/2)=0$.
Integrating the variance bound up to $1$ gives $(\log\zerLL)'(1)\le1/(2c)$, and
\eqref{zer:eq:Zformula} gives $Z\le1/c$ under the Riemann hypothesis for $L(s,\rho)$. Take
$\tau=4$ in \eqref{zer:eq:liminf}.
\end{proof}

\begin{theorem}[The cubic race; \stR]\label{zer:thm:cubicrace}
For the cubic weight $b(p)=\#\set{y\in\F_p:y^3-y-1=0}-1$, $p\ne23$, the following are
equivalent:
\begin{enumerate}[label=\textup{(\roman*)}]
\item every zero of $L(s,\rho)$ in the critical strip has real part $1/2$;
\item $J_4(x)\ge0$ for every sufficiently large real $x$;
\item $\displaystyle\liminf_{x\to\infty}\sqrt x\,J_4(x)\ge1-\frac9{\pi\sqrt{23}}>0.4026$.
\end{enumerate}
The same equivalence holds with strict positivity in \textup{(ii)}. Without any hypothesis,
the cubic $J_4$ takes positive values at arbitrarily large $x$.
\end{theorem}
\begin{proof}
The implication (i)$\Rightarrow$(iii) is \cref{zer:cor:cubicZ}, and (iii)$\Rightarrow$(ii) is
immediate, with strict positivity. For (ii)$\Rightarrow$(i), the Mellin transform of
$J_\tau$ is $(1-\tau^{s-1})/(s(s-1))$ times $\sum_{p\ne23}b(p)(\log p)p^{-s}$ for
$\Re s>1$; it has a pole at every zero of $L(s,\rho)$ with $1/2<\Re s<1$, it is regular on
the real interval $(1/2,\infty)$ because $\zerLL$ is positive there, and the Landau step
given in \cref{sec:hec} (an instance of \cref{ana:thm:onesided}) then shows that an
eventually nonnegative $J_\tau$ forces (i). An eventually nonpositive $J_\tau$ forces (i)
in the same way, by applying the argument to $-J_\tau$; for $\tau=4$ this would contradict
(iii), which proves the last assertion.
\end{proof}

The Riemann hypothesis for $L(s,\rho)$ is not a consequence of \cref{src:thm:dgrh}, and
establishing the sign in (ii) is the open step. The obstruction to replacing the cubic
weight by congruence weights is treated in
\cref{sec:hec}. \Cref{zer:prop:eta} gives one formula for all positive eta products; earlier
computations reach much farther for this particular field \cite{Ng2000Limiting}. The
quantity $2\zerLL'(1)/\zerLL(1)$ has also been enclosed numerically: with
$G(t)=f(it/\sqrt{23})$ and $G(1/t)=tG(t)$,
\[
 \zerLL(s)=\int_1^\infty G(t)\bigl(t^{s-1}+t^{-s}\bigr)\,\dd t,
\]
so $2\zerLL'(1)/\zerLL(1)$ is twice the ratio of the integrals of $G(t)\log t\,(1-1/t)$
and $G(t)(1+1/t)$, and \artifact{grh-transfer/checks/cubic23.py} encloses it in
$(0.32770,0.35161)$, consistently with the bound $1/c=0.39823\ldots$. The enclosure
evaluates the positive eta product on intervals of width $1/64$ up to $16$, using the first
$12$ factors of each infinite product; for omitted factors $1-u_i$ their product lies
between $1-\sum u_i$ and $1$, and for the integral tails one uses
$0<G(t)\le e^{-2\pi t/\sqrt{23}}$ and $\log t\le t$. These integrals and enclosures hold
independently of any hypothesis on zeros.

\subsection{Two further cubic tests}\label{hec:ss:othercubic}
{\emergencystretch=\textwidth\sources{grh-lean-transfers; mcnew_counterexample_and_grh_transfers}}

\paragraph{A binary event isolating one twisted factor.}
Let $H$ be the splitting field of $P(X)=X^3-X-1$ as before and $E=H(i)$, so that
$G=\Gal(E/\Q)=S_3\times C_2$. Let $\sigma$ be the sign character of $S_3$, $\rho$ its standard character and
$\eta$ the nontrivial character of $C_2$. For a prime $p\nmid46$ put $f(p)=1$ if either
$p\equiv1\pmod4$ and $P$ is irreducible modulo $p$, or $p\equiv3\pmod4$ and $P$ has three distinct linear
factors modulo $p$; put $f(p)=0$ otherwise. Evaluating on the three cycle types for each value of $\eta$
shows that the corresponding class function is
\begin{equation}\label{hec:eq:twistid}
 f=\frac{1+\sigma}4+\frac{(1+\sigma)\eta}{12}-\frac{\rho\eta}3 .
\end{equation}
Its mean is $1/4$, and it has no $\rho$ term.

\begin{theorem}[One Hecke $L$-function from one binary event; \stR]\label{hec:thm:twistedcubic}
Assume \cref{src:thm:dgrh}. All nontrivial zeros of $L(s,\rho\eta)$ lie on $\Re s=1/2$ if and only if for
every $\eps>0$ there is a constant $C_\eps$ such that, for $x\ge2$,
\begin{equation}\label{hec:eq:twistbound}
 \sum_{p\le x,\ p\nmid46}f(p)\log p-\frac x4\le C_\eps x^{1/2+\eps}.
\end{equation}
\end{theorem}

\begin{proof}
The standard representation of $S_3$ is induced from either nontrivial character of $A_3$, and twisting by
$\eta$ commutes with induction. Thus $L(s,\rho\eta)$ is the Hecke $L$-function of a nontrivial finite-order
character of $\Q(\sqrt{-23})$; it is entire, self-dual, and has the usual functional equation. Let
$\Psi_f(x)=\sum_{p^k\le x,\ p\nmid46}f(g_p^k)\log p$; the prime powers with $k\ge2$ contribute
$O(\sqrt x\log^2(2x))$. For $\Re s>1$ the transform of $\Psi_f(x)-x/4$ against $x^{-s-1}\,\dd x$ is
\[
 T(s)=-\frac1s\sum_\chi c_\chi\frac{L'}L(s,\chi)-\frac1{4(s-1)}+U(s),
\]
where the $c_\chi$ are the coefficients in \eqref{hec:eq:twistid} and $U$, which comes from the omitted
Euler factors, is holomorphic in $\Re s>0$. The pole at $1$ cancels, \cref{src:thm:dgrh} controls the four
linear terms, and at a zero $z$ of $L(s,\rho\eta)$ with $\Re z>1/2$, of multiplicity $m_z$, the residue of
$T$ is $m_z/(3z)$. Given \eqref{hec:eq:twistbound}, \cref{hec:lem:landau} shows that $T$ is holomorphic in
$\Re s>\theta$ for every $\theta>1/2$; the functional equation completes the forward direction. Conversely,
the explicit formula under RH for these finitely many Dirichlet and Hecke $L$-functions bounds their
centered sums by $O(\sqrt x\log^2(2x))$, and \eqref{hec:eq:twistid} gives \eqref{hec:eq:twistbound}.
\end{proof}

By the classification of admissible binary Frobenius events in \cref{sec:nab}, among such events this $f$
is the unique inversion-invariant $\set{0,1}$-valued class function of least support that has a negative
coefficient on $\rho\eta$ and coefficient zero on $\rho$; its support has three of the twelve elements of
$G$. The
identity and a three-cycle of $S_3$ have the same image in the abelianization of $G$ and different values
of $f$, which shows where an estimate beyond congruence conditions is needed.

\subsection{Dihedral conductors \texorpdfstring{$47$ and $71$}{47 and 71}}\label{hec:ss:dihedral}
{\emergencystretch=\textwidth\sources{grh_dihedral_47_71}}

At conductor $23$ there is one nonabelian Artin factor and its completion is positive on the real axis. At
conductors $47$ and $71$ there are two and three such factors. This subsection determines the real zeros of
their real linear combinations and shows that one ideal class still tests all factors at once.

\paragraph{Conductor $47$: the real zeros of the pencil.}
Let $K=\Q(\sqrt{-47})$ and let $H$ be its Hilbert class field. Reduction gives the forms
$(1,1,12)$, $(2,\pm1,6)$, $(3,\pm1,4)$, so the class group is cyclic of order five, with representatives up
to inversion
\[
 Q_0=x^2+xy+12y^2,\qquad Q_1=2x^2+xy+6y^2,\qquad Q_2=3x^2+xy+4y^2 .
\]
Complex conjugation acts by inversion, and $\Gal(H/\Q)$ is dihedral of order ten. Put
$B_i(z)=\sum_{x,y\in\Z}q^{Q_i(x,y)}$, $q=e^{2\pi iz}$, and let $\lambda=(1+\sqrt5)/2$,
$\mu=(1-\sqrt5)/2$. The two normalized eigenforms and their difference are
\begin{align}
 f_1&=(B_0-\mu B_1-\lambda B_2)/2, &
 f_2&=(B_0-\lambda B_1-\mu B_2)/2,\nonumber\\
 E(z)&=\eta(z)\eta(47z)=(B_1-B_2)/2=(f_1-f_2)/\sqrt5 \label{hec:eq:eta47}
\end{align}
\cite{BerkovichPatane2013Binary,BuzzardWeightOne}. The $L$-functions $L_1,L_2$ of $f_1,f_2$ are the Hecke
$L$-functions of the two inverse pairs of nontrivial class characters, that is, the two irreducible
nonabelian Artin factors of $H/\Q$. They have the same gamma factor, and we put
\[
 \Lambda_j(s)=\Bigl(\frac{\sqrt{47}}{2\pi}\Bigr)^{s}\Gamma(s)L_j(s),\qquad
 F_c(s)=\Lambda_1(s)-c\Lambda_2(s)\quad(c\in\R).
\]
These completions are entire. The one direction of the pencil omitted by this notation, a nonzero
multiple of $\Lambda_2$, has no real zero.

\begin{theorem}[Real zeros at conductor $47$; \stU]\label{hec:thm:real47}
For $m\ge0$ let $C_m=\Lambda_1^{(2m)}(1/2)/\Lambda_2^{(2m)}(1/2)$. The numbers $C_m$ decrease strictly to
$1$. The ratio
\[
 R_m(t)=\frac{\Lambda_1^{(2m)}(1/2+t)}{\Lambda_2^{(2m)}(1/2+t)}
\]
is positive and even, decreases strictly from $C_m$ to $1$ as $t$ increases from $0$ to $\infty$, and
satisfies $R_m''(0)<0$. Consequently the real zeros of $F_c^{(2m)}$ are as follows.
\begin{center}\begin{tabular}{ll}\toprule
Parameter & Real zeros with multiplicities \\\midrule
$c\le1$ or $c>C_m$ & none\\
$c=C_m$ & one double zero at $s=1/2$\\
$1<c<C_m$ & two simple zeros $s=1/2\pm t_c$, $t_c>0$\\\bottomrule
\end{tabular}\end{center}
For all real $s$, $\Lambda_1^{(2m)}(s)>\Lambda_2^{(2m)}(s)>0$. For odd orders, the ratio of the derivatives
of order $2m+1$ decreases strictly from $C_{m+1}$ to $1$ on $t>0$; the central zero of $F_c^{(2m+1)}$ is
simple unless $c=C_{m+1}$, when it is exactly triple, and there is one further simple symmetric pair of
real zeros exactly when $1<c<C_{m+1}$.
\end{theorem}

\begin{proof}
\emph{Positivity of the smaller kernel.}
For $y\ge1/\sqrt{47}$ one has $q=e^{-2\pi y}<r=41/100$. The first terms of $f_2$ are
$q-\lambda q^2-\mu q^3+\lambda q^4$, and the $n$th coefficient is bounded in absolute value by the number
of ideals of norm $n$ in $K$, hence by $d(n)\le n$. Dropping the positive third and fourth terms,
\[
 \frac{f_2(iy)}q\ge1-\tfrac53r-\frac{r^4(5-4r)}{(1-r)^2}>0.04 .
\]
The eta product is positive on the imaginary axis, so $f_1(iy)>f_2(iy)$ by \eqref{hec:eq:eta47}. Poisson
summation on each binary lattice gives $B_i(i/(47y))=\sqrt{47}\,yB_i(iy)$, because the Gram matrix has
determinant $47$ and its adjugate is integrally equivalent to it by a rotation. Thus positivity and the
ordering extend to all $y>0$, and both completions are positive at every real point.

\emph{A modular function with one pole.}
Put $h(z)=f_2(z)/E(z)$. Weights and characters cancel, and the Fricke transformations of numerator and
denominator agree, so $h$ descends to the compact curve $X_0(47)^+=X_0(47)/\langle W_{47}\rangle$. The eta
function has no zero in the upper half-plane; at infinity the numerator begins with $q$ and the
denominator with $q^2$. Since $47$ is prime, $X_0(47)$ has two cusps, exchanged by the Fricke involution;
their common image is unramified in the quotient and is the only pole of $h$, of order one. Hence $h$ is an
isomorphism of $X_0(47)^+$ with the sphere.

No elliptic point lies on the ray $iy$, $y>1/\sqrt{47}$. A matrix of determinant one fixing $iy$ has equal
diagonal entries, and a nonidentity elliptic element has absolute trace less than two; in $\Gamma_0(47)$
this would force both diagonal entries to vanish, which is incompatible with a lower-left entry divisible
by $47$ and determinant one. In the other coset write the matrix as
\[
 \begin{pmatrix}a&b\\47d&e\end{pmatrix}
 \begin{pmatrix}0&-1/\sqrt{47}\\\sqrt{47}&0\end{pmatrix}.
\]
Equality of the diagonal entries gives $b=-d$, and the trace bound forces $b=d=0$, whose only fixed point
is $i/\sqrt{47}$. The quotient map is therefore unramified along the open ray, so $h(iy)$ has nonzero
derivative there. It is real and tends to $+\infty$ like $q^{-1}$, so it increases strictly, and
\[
 b(u)=\frac{f_1(ie^u/\sqrt{47})}{f_2(ie^u/\sqrt{47})}=1+\frac{\sqrt5}{h(ie^u/\sqrt{47})}
\]
is strictly decreasing on $u>0$ and tends to one.

\emph{The ratio under the Mellin integral.}
Let $K_j(u)=e^{u/2}f_j(ie^u/\sqrt{47})$. By the Fricke identities these positive functions are even, and
their double-exponential decay justifies all differentiations in
\[
 \Lambda_j(1/2+t)=\int_\R K_j(u)e^{tu}\,\dd u .
\]
Thus $R_m(t)$ is the expectation of $b(u)$ under the probability density proportional to
$u^{2m}K_2(u)\cosh(tu)$ on $u>0$, and its derivative is the covariance of $b(u)$ with $u\tanh(tu)$. For
$t>0$ the second function increases strictly and $b$ decreases strictly; integrating
$(b(u)-b(v))(u\tanh(tu)-v\tanh(tv))/2$ over two independent draws shows that the covariance is negative.
At $t=0$ the same calculation after a second differentiation gives $R_m''(0)=\hecCov(b(u),u^2)<0$, and
weighting the measure at $t=0$ by $u^2$ proves $C_{m+1}<C_m$. As $t\to\infty$ with $m$ fixed, or
$m\to\infty$ with $t=0$, these probability measures leave every compact interval: the ratio of the mass on
$[0,A]$ to the mass on $[A+1,A+2]$ tends to zero by the exponential tilt $\cosh(tu)$, respectively the
power tilt $u^{2m}$. Since $b$ is bounded, continuous and tends to one, both limits are one. The denominator
$\Lambda_2^{(2m)}$ is positive on the real line, so the strict derivative signs away from the centre and the
strict second derivative at the centre give the asserted multiplicities.

\emph{Odd derivatives.}
For $t>0$ the ratio of the derivatives of order $2m+1$ is the expectation of $b(u)$ under the density
proportional to $u^{2m+1}K_2(u)\sinh(tu)$, and its derivative is the covariance with $u\coth(tu)$, a
strictly increasing function of $u$ because $\cosh v\sinh v>v$ for $v>0$. The ratio decreases strictly from
$C_{m+1}$ to $1$, by dominated convergence at zero and the same escape argument at infinity. The central
zero of $F_c^{(2m+1)}$ is simple unless $c=C_{m+1}$; in that case its cubic Taylor coefficient is nonzero
because $C_{m+2}<C_{m+1}$.
\end{proof}

The transition value for the functions themselves is certified by rational arithmetic. With
$h_0=2\pi/\sqrt{47}$ and $G(a)=\int_1^\infty e^{-av^2}\,\dd v$, the substitution $v=e^{u/2}$ gives
$\Lambda_j(1/2)=4\sum_{n\ge1}a_j(n)G(h_0n)$, where $a_j(n)$ are the coefficients of $f_j$. Since
$v^2\ge1+2(v-1)$ for $v\ge1$, $G(h_0n)\le e^{-h_0n}/(2h_0n)$, and with $\abs{a_j(n)}\le n$ and
$e^{-h_0}<r=41/100$ the tail after $N=28$ terms, before multiplication by $4$, is at most
$r^{N+1}/(2h_0(1-r))$. The value $G(a)$ is obtained by subtracting $\int_0^1e^{-av^2}\,\dd v$ from
$\tfrac12\sqrt{\pi/a}$, with the integrated Taylor polynomials $\sum_{k=0}^{M}(-a)^k/(k!(2k+1))$ as upper
bound at $M=122$ and lower bound at $M=123$. The resulting enclosures are
\begin{equation}\label{hec:eq:C0}
\begin{gathered}
 0.6660152520\le\Lambda_1(1/2)\le0.6660152521,\qquad
 0.4586804524\le\Lambda_2(1/2)\le0.4586804526,\\
 1.4520244940\le C_0\le1.4520244943
\end{gathered}
\end{equation}
(\artifact{grh_dihedral_47_71/check_threshold.py}). For example, $F_{7/5}$ has two simple real zeros,
while $F_{3/2}$ has none.
{\emergencystretch=14em\par}

\begin{remark}\label{hec:rem:bookerthorne}
Positivity of a Mellin kernel does not control zeros off the real axis. For $c=1$ the completion $F_1$ is,
up to the factor $\sqrt5$, the completed $L$-function of $E=\eta(z)\eta(47z)$, and it is positive on the
real axis. But $E$ is a combination of two Hecke eigenforms with different $T_2$-eigenvalues, so by the
theorem of Booker and Thorne \cite[Theorem~1.1]{BookerThorne2014Zeros} its Dirichlet series has zeros with
$\Re s>1$; each of them recurs in disjoint translates of a small disk by almost periodicity and Rouch\'e's
theorem. These zeros belong to the difference of the factors.
\end{remark}

\paragraph{Conductor $47$: one quadratic form tests both factors.}
Choose $Q$ to be $Q_1$ or $Q_2$, once and for all, and let
\[
 S_Q(x)=\sum_{\substack{p\le x,\ p\ne47\\ p=Q(u,v)\ \text{for some }u,v\in\Z}}\Bigl(1-\frac px\Bigr)\log p,
\]
each prime being counted once. The corresponding conjugacy class has two of the ten group elements, so the
density is $1/5$, and the linear smoothing halves the main term.

\begin{theorem}[One class, two factors; \stR]\label{hec:thm:primes47}
Assume \cref{src:thm:dgrh}. For either choice of $Q$ the following are equivalent.
\begin{enumerate}[label=\textup{(\arabic*)}]
\item The nontrivial zeros of both $L_1$ and $L_2$ lie on $\Re s=1/2$.
\item For every $\eps>0$ there is $A_\eps$ such that $S_Q(x)\le x/10+A_\eps x^{1/2+\eps}$ for all large
$x$.
\item For every $\eps>0$ there is $A_\eps$ such that $S_Q(x)\ge x/10-A_\eps x^{1/2+\eps}$ for all large
$x$.
\end{enumerate}
In particular either direction of the error suffices by itself, and common zeros of the two factors cannot
defeat the test.
\end{theorem}

\begin{proof}
Let $r$ generate the rotations of the dihedral group, labelled so that the class of $Q_1$ corresponds to
$r^{\pm1}$. By class field theory the unramified rational primes represented by $Q$ are exactly those with
Frobenius in $C=\set{r,r^{-1}}$, respectively $C=\set{r^2,r^{-2}}$. Put $a=(\sqrt5-1)/2$ and
$b=-(\sqrt5+1)/2$. For the first class, column orthogonality gives
\[
 5\,\hecind_C=\triv+\hecsgn+a\chi_1+b\chi_2,
\]
where $\chi_j$ is the Artin character of $f_j$; for the second class $a$ and $b$ are interchanged. Here
$a+b=-1$, and $a,b$ are linearly independent over $\Q$.

Let $U(x)$ be the sum of $\hecind_C(g_p^k)(1-p^k/x)\log p$ over $p^k\le x$, $p\ne47$; Chebyshev's estimate
gives $U-S_Q=O(\sqrt x)$. For $\Re s>1$,
\begin{align*}
 \int_1^\infty U(x)x^{-s-1}\,\dd x&=\frac{D(s)}{s(s+1)},\\
 5D(s)&=-\frac{\zeta'}\zeta(s)-\frac{L'}L(s,\chi_{-47})-a\frac{L_1'}{L_1}(s)-b\frac{L_2'}{L_2}(s)+B(s),
\end{align*}
where the ramified correction $B$ is holomorphic in $\Re s>0$. The residue at $s=1$ is $1/10$; subtract
$1/(10(s-1))$ and call the result $M(s)$. By \cref{src:thm:dgrh}, \cref{hec:thm:real47} (which gives
$\Lambda_j>0$ on the real axis) and the cancellation at $1$, the function $M$ is regular at every real
point of $(1/2,\infty)$. At a nonreal $\rho$ with $\Re\rho>1/2$ its residue is
\[
 -\frac{a\ord_\rho L_1+b\ord_\rho L_2}{5\rho(\rho+1)},
\]
and the numerator vanishes only if both integer multiplicities vanish. Take $\sigma_0=1/2+\eps<1$. Under
(2) or (3), \cref{hec:lem:landau} applies to $U(x)-x/10$ or to its negative, with no real pole to consider,
and shows that $M$ is holomorphic in $\Re s>\sigma_0$. The residue formula excludes all zeros there; let
$\eps\to0$ and use the functional equations. Conversely $\zeta_H=\zeta L(s,\chi_{-47})L_1^2L_2^2$, and
under (1) and \cref{src:thm:dgrh} the effective Chebotarev theorem for this fixed field
\cite{GrenieMolteni2017Effective} gives the weighted count $x/5+O(\sqrt x\log^2(2x))$; partial integration
against $1-p/x$ gives both one-sided estimates.
\end{proof}

No rational congruence weight reproduces this class, even in mean square. Let $m$ be any modulus and $v$
any function of a reduced residue class modulo $m$; since $v$ is also a function of the class modulo
$\lcm(47,m)$, we may assume $47\mid m$. The intersection of $H$ with $\Q(\zeta_m)$ is then $K$, the maximal
abelian subextension of $H/\Q$. Chebotarev's theorem in the compositum gives
the uniform distribution on the five rotations for residue classes that split in $K$; the conditional
probability of $C$ is $2/5$ there and $0$ on inert classes. Therefore
\begin{equation}\label{hec:eq:325}
 \lim_{X\to\infty}\frac1{\pi(X)}\sum_{\substack{p\le X\\p\nmid47m}}\abs{\hecind_C(g_p)-v(p\bmod m)}^2
 \ge\frac12\cdot\frac25\cdot\frac35=\frac3{25},
\end{equation}
with equality for $m=47$ and the conditional mean. An exhaustive enumeration through $10^6$ finds $7{,}819$,
$15{,}692$ and $15{,}690$ primes represented by $Q_0,Q_1,Q_2$, with $47$ omitted, and the centered values
$(S_Q(x)-x/10)/\sqrt x$ for $Q_1,Q_2$ are about $-0.04823$ and $-0.11582$ at this cutoff
(\artifact{grh_dihedral_47_71/sample_primes.py}). These counts check the identification of the classes;
the theorem requires a bound at all large $x$.
{\emergencystretch=14em\par}

\paragraph{A real-zero bound at prime levels.}
Let $p\equiv23\pmod{24}$ be prime and $m=(p+1)/24$. Let $F$ be a nonzero cusp form with real coefficients,
of weight one on $\Gamma_0(p)$ with character $\legendre{-p}{\cdot}$, satisfying the Fricke law
$F(i/(py))=\sqrt p\,yF(iy)$ for $y>0$. This covers the theta eigenforms of nonreal class characters of
$\Q(\sqrt{-p})$. Let $\Lambda_F(s)=(\sqrt p/(2\pi))^s\Gamma(s)L(s,F)$.

\begin{theorem}[Real zeros and interpolation; \stU]\label{hec:thm:interp}
For each integer $k\ge0$, the entire function $x\mapsto\Lambda_F^{(2k)}(1/2+\sqrt x)$ has at most $m-1$
zeros on $[0,\infty)$, counted with multiplicity. Equivalently, $\Lambda_F^{(2k)}$ has at most $2(m-1)$
zeros on the real axis, counted with multiplicity, a zero of order $r$ at $x=0$ corresponding to a central
zero of order $2r$.

At $p=71$, prescribe any two points of $[0,\infty)$ in the $x$-coordinate, repetition allowed. There is a
real linear combination of the three completed nonabelian class-character factors, unique up to a nonzero
scalar, for which $x\mapsto\Lambda_F^{(2k)}(1/2+\sqrt x)$ has exactly that multiset of zeros.
\end{theorem}

\begin{proof}
\emph{Divisor bound.}
The criterion for eta quotients, applied to $E(z)=\eta(z)\eta(pz)$, uses $p+1\equiv0\pmod{24}$ in both
congruences and gives the weight, the character and the Fricke law of $F$. The quotient $F/E$ descends to
the Fricke quotient $X_0(p)^+$. The product $E$ does not vanish in the upper half-plane and begins with
$q^m$. The two cusps of $X_0(p)$ are exchanged by the Fricke involution, so $F/E$ has a single possible
pole on the quotient, of order at most $m-1$, and at most that many zeros if it is not constant.

The arc $iy$, $y>1/\sqrt p$, injects into the quotient. Indeed, for
$\left(\begin{smallmatrix}a&b\\c&d\end{smallmatrix}\right)\in\Gamma_0(p)$ a purely imaginary image of $iy$
requires $acy^2+bd=0$. If both terms are nonzero, then $ad$ and $bc$ are nonzero of opposite signs, which
is incompatible with $ad-bc=1$. Otherwise both products vanish and the matrix is diagonal or antidiagonal;
the first case is the identity in $\PSL_2$ and the second is impossible since $p\mid c$. In the other
Fricke coset one first replaces $y$ by $1/(py)$, below the end of the arc. Thus $F(iy)$ has at most $m-1$
sign changes on the open arc.

\emph{Passage to the Mellin integral.}
For $u=\log(\sqrt p\,y)>0$ put $K(u)=e^{u/2}F(ie^u/\sqrt p)$. The centred derivative is
$2\int_0^\infty u^{2k}K(u)\cosh(tu)\,\dd u$, and the cusp expansion gives double-exponential decay. For
$0\le t_1<\dots<t_l$ and $0<u_1<\dots<u_l$,
\[
 \det\bigl[\cosh(t_iu_j)\bigr]_{i,j=1}^l>0 .
\]
To see this, expand each entry in the terms $(t_i^2)^h(u_j^2)^h/(2h)!$, truncated at $h=M$. The
Cauchy--Binet formula writes the determinant as a sum of products of two generalized Vandermonde
determinants with positive weights; each product is nonnegative by the alternant formula for Schur
polynomials, and the term with $h=0,\dots,l-1$ is positive. This lower bound persists as $M\to\infty$. The
determinant remains positive for repeated nodes in $x_i=t_i^2$ when the rows at a node of multiplicity $r$
are replaced by $\partial_x^j\cosh(\sqrt xu)/j!$, $0\le j<r$: approximate each repeated node by a cluster
of distinct nonnegative $x$'s and divide by the product of differences within each cluster. Taylor
expansion in $x$ gives the derivative rows in the limit, including at $x=0$; the lower bound supplied by
$h=0,\dots,l-1$ contains the full Vandermonde determinant in the $x$'s, and after division the factors
between different clusters have positive limits.

Let $u_1<\dots<u_b$ be the sign changes of $K$, so $b\le m-1$. If the transform had $b+1$ zeros in $x$,
counted with multiplicity, take the corresponding rows of the confluent determinant, with columns at
$u_1,\dots,u_b,u$ and the last point variable. By the strict sign of ordered minors this determinant
changes sign exactly when $u$ crosses the $b$ points, so its product with $u^{2k}K(u)$ has constant sign
and nonzero integral. Expanding along the last column expresses this integral as a linear combination of
the vanishing Mellin integrals and their $x$-derivatives, a contradiction. The same argument excludes an
identically vanishing transform.

\emph{Sharpness at level $71$.}
The class group has order seven, and its six nontrivial characters form three inverse pairs. The three
theta eigenforms are linearly independent, because their $T_2$-eigenvalues are the distinct numbers
$2\cos(2j\pi/7)$, $1\le j\le3$ \cite{BerkovichPatane2013Binary}, and they satisfy the Fricke law by Poisson
summation. Two homogeneous linear conditions on three real coefficients have a nonzero solution; for a
repeated point the conditions are the vanishing of the value and of the first $x$-derivative. The bound
$m-1=2$ excludes any further zero or multiplicity. If the solution space had dimension two, imposing
vanishing at a third point would leave a nonzero solution with three zeros.
\end{proof}

The corresponding geometric and Mellin-transform method in weight at least two is due to Wang
\cite{Wang2026Modular}; here the eta divisor supplies the geometric bound in weight one, and the factor
$u^{2k}$ gives the same count for every even derivative.

\begin{example}[A central zero of order four]\label{hec:ex:fourth}
Order the three factors at level $71$ so that their $T_2$-eigenvalues are
$2\cos(6\pi/7),2\cos(2\pi/7),2\cos(4\pi/7)$, and let $A=(\Lambda_j(1/2))_j$ and
$B=(\Lambda_j''(1/2))_j$. By the interpolation theorem with $k=0$ at the repeated node $x=0$, the
conditions $w\cdot A=w\cdot B=0$ have a one-dimensional solution space, so $A\times B\ne0$, and
\[
 \mathcal L(s)=\det\begin{pmatrix}A_1&A_2&A_3\\ B_1&B_2&B_3\\
 \Lambda_1(s)&\Lambda_2(s)&\Lambda_3(s)\end{pmatrix}
\]
has vanishing value and second derivative at $s=1/2$; the odd central derivatives vanish by symmetry. The
zero count gives $\mathcal L^{(4)}(1/2)\ne0$ and excludes every other real zero of $\mathcal L$. The
central cancellation takes place between three distinct Artin factors. Numerically the coefficient vector
is approximately $(-0.3172807331,-0.6479315312,1)$ (\artifact{grh_dihedral_47_71/example71.py}); the proof
does not depend on these values.
{\emergencystretch=14em\par}
\end{example}

\paragraph{Conductor $71$: positivity and a one-class test.}
Let $H$ now be the Hilbert class field of $\Q(\sqrt{-71})$. A reduced form has $a\le\sqrt{71/3}<5$, and
the seven classes are represented, up to inversion, by
\[
 Q_0=(1,1,18),\quad Q_1=(2,1,9),\quad Q_2=(4,3,5),\quad Q_3=(3,1,6).
\]
The class group is cyclic of order seven. A prime above $2$ generates it; the square of that prime has norm
$4$ and nontrivial class, and the only nonprincipal inverse pair representing $4$ is $Q_2$, which matches
the indices with the powers of a generator up to inversion. The group $\Gal(H/\Q)$ is dihedral of order
$14$. With $\Theta_k(iy)=\sum_{x,v\in\Z}\exp(-2\pi yQ_k(x,v))$, the three weight-one Hecke forms are
\[
 f_j(iy)=\tfrac12\Theta_0(iy)+\sum_{k=1}^3\cos(2\pi jk/7)\,\Theta_k(iy),\qquad1\le j\le3.
\]

\begin{proposition}[Positivity at conductor $71$; \stU]\label{hec:prop:pos71}
With $q=e^{-2\pi y}$, for $0<q\le12/25$,
\[
 f_1(iy)/q\ge0.999,\qquad f_2(iy)/q\ge0.877,\qquad f_3(iy)/q\ge0.324 .
\]
Consequently each $f_j(iy)$ is positive for all $y>0$, and the three completed $L$-functions are positive
on the whole real axis.
\end{proposition}

\begin{proof}
For $y\ge1/\sqrt{71}$ one has $q<12/25$, since $2\pi/\sqrt{71}>37/50$ and the Taylor polynomial of degree
four of the exponential at $37/50$ exceeds $25/12$. Let $r_k(n)$ be the number of representations of $n$ by
$Q_k$ and $t=2\cos(2\pi j/7)$. Twice the $n$th coefficient of $f_j$ is
$r_0(n)+tr_1(n)+(t^2-2)r_2(n)+(1-t-t^2)r_3(n)$, and the three roots of $t^3+t^2-2t-1$ lie in
$[1.2469,1.2470]$, $[-0.4451,-0.4450]$ and $[-1.8020,-1.8019]$ by sign changes. The coefficients up to
$N=24$ are enumerated from $4aQ(x,v)=(2ax+bv)^2+71v^2$. All coefficients are bounded in absolute value by
the number of ideals of norm $n$, hence by $n$; with $r=12/25$ the remaining normalized series is at most
$r^N((N+1)/(1-r)+r/(1-r)^2)$. Interval Horner evaluation on the $24$ consecutive intervals of width $1/50$,
with this tail subtracted, gives the stated bounds (\artifact{grh_dihedral_47_71/check_level71.py}, in
rational arithmetic).{\emergencystretch=14em\par}
Poisson summation gives $\Theta_k(i/(71y))=\sqrt{71}\,y\Theta_k(iy)$, so positivity
extends to the whole imaginary axis, and the Mellin integral is positive at every real argument.
\end{proof}

For a dihedral group of order $2n$ with $n$ prime, a single nonidentity rotation class detects every
nonzero vector of nonnegative integral zero multiplicities, by Galois conjugation of the character values
(\cref{sec:nab}); when the nonlinear factors have no real zero in $(1/2,1)$, a one-sided bound suffices.
\Cref{hec:prop:pos71} supplies the real-axis hypothesis unconditionally. Hence, under \cref{src:thm:dgrh},
for any one of the three nonprincipal inverse pairs $[a]$ of classes the upper bound
\begin{equation}\label{hec:eq:71bound}
 \sum_{\substack{p\le x\\ \Frob_p\in[a]}}\log p\,\Bigl(1-\frac px\Bigr)\le\frac x{14}+O_\eps(x^{1/2+\eps})
 \qquad\text{for every }\eps>0
\end{equation}
is equivalent to RH for all three nonlinear factors of $\zeta_H$. The bound \eqref{hec:eq:71bound} is open.

\subsection{Eta products and Jensen polynomials}\label{hec:ss:jensen}
{\emergencystretch=\textwidth\sources{grh_transfers_density_gaps_and_zero_tests; grh_dihedral_47_71; mcnew_density_and_grh_transfers}}

For an even kernel $K$ on $\R$ with sufficient decay write
\begin{equation}\label{hec:eq:jensendef}
 \int_\R e^{zu}K(u)\,\dd u=\sum_{n\ge0}\frac{\gamma_n}{n!}z^{2n},\qquad
 J^{d,n}(X)=\sum_{j=0}^d\binom dj\gamma_{n+j}X^j .
\end{equation}
The polynomials $J^{d,n}$ are the Jensen polynomials of degree $d$ and shift $n$. If
$G(w)=\int_\R\cosh(\sqrt w\,u)K(u)\,\dd u$, then $\gamma_n=G^{(n)}(0)$. For the completed $L$-function
$\Lambda$ of \eqref{hec:eq:Lambda23} one has $G(w)=\Lambda(1/2+\sqrt w)$, a function of order at most $1/2$
whose zeros are negative real numbers exactly when the zeros of $\Lambda$ are on the critical line. By
Hadamard factorization and Jensen's criterion, RH for $L$ is therefore equivalent to real-rootedness of
$J^{d,n}$ for all $d$ and $n$. This subsection records what the kernel alone gives.

\paragraph{The quadratic inequalities.}
Let $h(v)=v/4+\log\eta(ie^v)$. The inversion formula for $\eta$ shows that $h$ is even, and for $v\ge0$,
with $Y=2\pi e^v$,
\begin{equation}\label{hec:eq:hseries}
 h(v)=\frac v4-\frac Y{24}-\sum_{m\ge1}\sigma_{-1}(m)e^{-mY},\qquad\sigma_{-1}(m)=\sum_{d\mid m}d^{-1}.
\end{equation}

\begin{theorem}[Quadratic Jensen inequalities for paired eta kernels; \stU]\label{hec:thm:turan}
Let $g(u)=\sum_jr_j\bigl[h(u-a_j)+h(u+a_j)\bigr]$ for a nonempty finite list of $a_j\in\R$ and $r_j>0$, and
let $K=e^{g}$. Then all $\gamma_n$ are positive and $\gamma_{n+1}^2>\gamma_n\gamma_{n+2}$ for every
$n\ge0$; equivalently, $J^{2,n}$ has two distinct negative real roots for every $n$. In particular this
holds for the centred Mellin kernel of every product $\eta(az)\eta(bz)$ with positive integers $a\le b$,
for which $g(u)=h(u-v)+h(u+v)$ with $v=\tfrac12\log(b/a)$, and so for the coefficients of
$\Lambda(1/2+z)$ in \eqref{hec:eq:Lambda23}.
\end{theorem}

\begin{proof}
The series \eqref{hec:eq:hseries} and its first four derivatives converge locally uniformly. The fourth
derivative in $v$ of $e^{-mY}$ is $e^{-q}q\bigl((q-6)q^2+7q-1\bigr)$ with $q=mY\ge2\pi>6$, hence positive.
Consequently $h^{(4)}(v)<0$ for all real $v$, by evenness for $v<0$. Then $g$ is even with $g^{(4)}<0$, so
$g'''(u)<0$ for $u>0$ and $(ug''-g')'=ug'''<0$. Since $g'(0)=0$,
\[
 \bigl(\log K(\sqrt x)\bigr)''=\frac{ug''(u)-g'(u)}{4u^3}<0,\qquad u=\sqrt x>0 .
\]
At both ends $K$ decays at least as fast as $\exp(-ce^{\abs u})$ for some $c>0$. This is the eta-kernel
curvature argument of Conrey and Ghosh \cite{ConreyGhosh1994Turan}; the symmetric translations preserve
the sign of the fourth derivative for unequal dilations.

Write $V(x)=-\log K(\sqrt x)$ and $M(p)=\int_0^\infty x^{p-1}e^{-V(x)}\,\dd x$ for $p>0$, and let $E_p$
denote expectation for the probability density $x^{p-1}e^{-V(x)}/M(p)$. Integration by parts gives
$E_{p+1}[V']=pM(p)/M(p+1)$. Since $V'$ is strictly increasing,
\[
 E_{p+2}[V']-E_{p+1}[V']=\frac{\hecCov_{p+1}(x,V'(x))}{E_{p+1}[x]}>0,
\]
as one sees by averaging $(x-y)(V'(x)-V'(y))/2$ over two independent draws. Hence the consecutive ratios
of $M(p)/\Gamma(p)$ decrease strictly as $p$ increases by one. Since
\[
 \gamma_n=\frac{n!}{(2n)!}M(n+1/2)=\frac{\sqrt\pi}{4^n}\frac{M(n+1/2)}{\Gamma(n+1/2)},
\]
this gives $\gamma_{n+1}^2>\gamma_n\gamma_{n+2}$, and all $\gamma_n$ are positive.
\end{proof}

The same calculation applies to every positive eta product whose exponents are paired under
$d\leftrightarrow N/d$, after centring its Mellin kernel at half the weight. Such a kernel need not come
from an eigenform, and the next result shows that the quadratic inequalities carry no information about
zeros off the real axis.

\begin{theorem}[Two-factor eta products; \stU]\label{hec:thm:etafamily}
For positive integers $a\le b$ let $f_{a,b}(z)=\eta(az)\eta(bz)$ and
$\Lambda_{a,b}(s)=(\sqrt{ab})^s\int_0^\infty f_{a,b}(iy)y^{s-1}\,\dd y$. The function $\Lambda_{a,b}$ is
zero-free in $\Re s>1$ if and only if
\[
 b/a\in\set{1,2,3,5,7,11,23}.
\]
For every other ratio it has infinitely many zeros in $\Re s>1$, and at every fixed shift $n$ the number
of nonreal roots of $J^{d,n}$ tends to infinity with $d$; in fact it is $\gg_{a,b,n}\log d/\log\log d$.
In the coprime subfamily with $24\mid a+b$, the only surviving pair is $(1,23)$.
\end{theorem}

\begin{proof}
A common dilation of $a$ and $b$ leaves the normalized completion unchanged, so divide by $\gcd(a,b)$. Set
$g_0=\gcd(a+b,24)$, $r=24/g_0$, $A=ra$, $B=rb$, $N=AB$ and $m=(a+b)/g_0$. Both eta congruence sums for
$f_{A,B}$ at level $N$ equal $A+B=24m$, and the cusp-order formula has positive summands, so $f_{A,B}$ is a
weight-one cusp form with character $\legendre{-AB}{\cdot}$ \cite[Theorems~1.64--1.65]{Ono2004Web}. Its
expansion begins with $q^m$ with coefficient one; write $a_n$ for its coefficients. One has
$\gcd(m,N)=1$, because $a+b$ is coprime to $ab$ and $m$ is coprime to $r$. If $m>1$, the coefficient of $q$
in $T_mf_{A,B}$ is $a_m=1$ whereas $a_1=0$, so $f_{A,B}$ is not an eigenfunction of all $T_p$ with
$p\nmid N$, which generate $T_m$.

If $m=1$, then $(A,B)$ is one of the twelve partitions of $24$ into two parts and $a_1=1$. The four pairs
$(5,19),(7,17),(9,15),(11,13)$ have $a_2=a_4=0$, and the pair $(10,14)$ has $a_3=a_9=0$; the prime $p$
involved is coprime to $N$, so the eigenform relation $a_{p^2}=a_p^2-\chi(p)$ fails. The remaining seven
pairs are the classical multiplicative eta products of weight one
\cite{DummitKisilevskyMcKay1985Multiplicative,VoskresenskayaEta}; their ratios form the displayed set, and
the absolutely convergent Euler products give nonvanishing in $\Re s>1$. Each excluded form fails to be
a Hecke eigenform, and by \cite[Theorem~1.1 and Remark~(2)]{BookerThorne2014Zeros} its $L$-function has
infinitely many zeros in $\Re s>1$.

For the Jensen polynomials of an excluded pair, write the completed transform as $\mathcal E(s)L(s)$ with
$\mathcal E(s)=(\sqrt N/2\pi)^s\Gamma(s)$ and $L$ the Dirichlet series of $f_{A,B}$. Choose a closed disk
$D\subset\set{\Re s>1}$ with a zero of $L$ inside and none on its boundary. Absolute convergence gives
arbitrarily large $t_j$ with $L(s+it_j)\to L(s)$ uniformly on $D$: truncate the series and approximate
the phases $n^{-it_j}$ by one for the finitely many retained indices. Every fixed derivative of $L$ is
bounded on the corresponding vertical strip. Put $w=(s-1/2)^2$ and fix $n\ge0$. Uniformly on $D+it_j$ as
$j\to\infty$,
\[
 \frac{[2(s-1/2)]^nG^{(n)}(w)}{\mathcal E(s)(\log s)^n}=L(s)+O_n(1/\log t_j).
\]
Indeed, expanding $\dd/\dd w=[2(s-1/2)]^{-1}\dd/\dd s$ repeatedly, the leading term after multiplication
by the displayed power is $(\mathcal EL)^{(n)}$ and the terms with $k<n$ derivatives have coefficients
$O_n(\abs s^{k-n})$; Stirling's formula gives
$\mathcal E^{(k)}/\mathcal E=(\log s)^k+O_k((\log\abs s)^{k-1})$ on this strip, and Leibniz's rule
completes the estimate. Rouch\'e's theorem gives a zero of $G^{(n)}((s-1/2)^2)$ in every translated disk
of large height. Disjoint disks of increasing positive ordinate have squares about $1/2$ off the real axis,
so $G^{(n)}$ has infinitely many nonreal zeros. Finally $J^{d,n}(w/d)\to G^{(n)}(w)$ locally uniformly as
$d\to\infty$, by termwise binomial scaling and domination by the absolute Taylor series, so any finite
collection of disjoint small disks around nonreal zeros contains zeros of every scaled polynomial of
sufficiently large degree.

For the quantitative bound, a fixed uniform approximation on $D$ suffices for Rouch\'e's theorem, and the
corresponding shifts are relatively dense: a finite Dirichlet sum approximates $L$ uniformly, and the
returns of its phase flow to a neighbourhood of the identity have bounded gaps, by compactness of the
closure of the flow in a finite-dimensional torus. One can therefore choose $\gg T$ disjoint translated
disks at heights between $T/2$ and $T$, on whose boundaries $\abs{G^{(n)}((s-1/2)^2)}\ge e^{-2T}$ for large
$T$. Their squares lie in $\abs w\le R$ with $R=O(T^2)$. The binomial scaling coefficients
$u_{d,j}=\prod_{k=0}^{j-1}(1-k/d)$ satisfy $0\le1-u_{d,j}\le j(j-1)/(2d)$, and positivity of the Taylor
coefficients gives
\[
 \max_{\abs w\le R}\abs{G^{(n)}(w)-J^{d,n}(w/d)}\le\frac{R^2G^{(n+2)}(R)}{2d}\le\frac{e^{C_nT\log T}}d .
\]
The last inequality uses the bound $C\exp(-ce^{\abs u})$ for the centred kernel, which gives
$\log\max_{\abs w\le R+1}\abs{G(w)}=O(\sqrt{R+1}\log(R+2))$, and Cauchy's bound for the derivatives. With
$T=c_n\log d/\log\log d$ and $c_n$ small, the approximation error is at most $d^{-1/2}<e^{-2T}$, and a
second application of Rouch\'e's theorem transfers all these disks to the polynomial of degree $d$.

In the coprime subfamily with $24\mid a+b$ one has $r=1$ and $(A,B)=(a,b)$, and among the seven surviving
partitions of $24$ only $(1,23)$ has coprime parts.
\end{proof}

\paragraph{Finite ranges of degrees at conductor $23$.}
Let $\Lambda$ be as in \eqref{hec:eq:Lambda23}, with centred kernel
\[
 K(u)=e^{u/2}\,\eta\bigl(ie^u/\sqrt{23}\bigr)\,\eta\bigl(i\sqrt{23}\,e^u\bigr).
\]
The location of the low zeros gives real-rootedness in a range of degrees, uniformly in the shift.

\begin{theorem}[Jensen polynomials through degree sixteen; \stU]\label{hec:thm:sixteen}
For the conductor-$23$ kernel, $J^{d,n}$ has only negative real zeros for all $1\le d\le16$ and all
$n\ge0$.
\end{theorem}

\begin{proof}
\emph{A zero-free rectangle.} Put $a=\Re s-1/2$ and $t=\Im s$. Then
\[
 \tfrac12\Re\Lambda(s)=\int_0^\infty K(u)\cosh(au)\cos(tu)\,\dd u,
\]
and on $\abs a\le1/2$, $\abs t\le4$ the right side exceeds $1/1000$. This is certified with integer bounds
(\artifact{grh_transfers_density_gaps_and_zero_tests/checks/check_height.py}) as follows. Subdivide
$0\le u\le4$ in steps of $1/1024$ and $0\le t\le4$ in steps of $1/256$. On each $u$-interval bound
$q=\exp[-(2\pi/\sqrt{23})e^u]$ using $1.3101<2\pi/\sqrt{23}<1.3102$. In
$K=e^{u/2}q\prod_{m\ge1}(1-q^m)(1-q^{23m})$ keep the first eight factors of the first product; the product
of the remaining factors lies between $1-2q^9/(1-q)$ and $1$. Bound the cosine from below by its lower
Taylor polynomial of degree $62$ at $\lfloor512t_{\min}u_{\min}\rfloor/512$, minus the maximal phase
distance from this point in the cell. A positive lower bound for the cosine is multiplied by the lower
bound for $K$, a negative one by the upper bound for $K\cosh(u/2)$. The sum of the contributions is a
lower integral for every $a$ and every $t$ in the cell. Beyond $u=4$ the integrand is bounded by
$e^ue^{-e^u}$, with tail $e^{-e^4}<2^{-64}$. All exponentials use a rational Taylor bound on $[0,1]$
followed by seven squarings with outward rounding.
{\emergencystretch=14em\par}

\emph{The sector argument.} The Euler product and the functional equation place the zeros of $\Lambda$ in
$\abs a\le1/2$, and the zeros of $G$ are the numbers $w=(a+it)^2$. If all zeros have $\abs t\ge T$, with
$T>1/2$, then $\abs a/\abs t\le1/(2T)$ and the zeros of $G$ lie in the closed sector about the negative
real axis with half-angle $\theta=2\arctan(1/(2T))$. The estimate
$\log\max_{\abs z\le R}\abs{\Lambda(1/2+z)}=O(R\log(2+R))$ shows that $G$ has order at most $1/2$, so it is
the locally uniform limit of its genus-zero partial products, which are real polynomials after grouping
conjugate zeros. The closed sector is convex, so by the Gauss--Lucas theorem the zeros of every $G^{(n)}$
lie in it, and all $G^{(n)}$ have positive Taylor coefficients. The sector theorem of Obreschkoff and
Chasse \cite{Chasse2013Laguerre} gives real-rooted Jensen polynomials for each $G^{(n)}$ through degree
\[
 d\le\csc^2\theta=\Bigl(T+\frac1{4T}\Bigr)^2 .
\]
With $T=4$ this is $d\le\lfloor(65/16)^2\rfloor=16$. Positivity of the coefficients puts the roots on the
negative axis.
\end{proof}

The same sector calculation extends the range when zero computations are taken as input. The ten low
ordinates $5.1156833,\dots,18.0806617$ listed by the LMFDB for the newform $23.1.b.a$, with the stated
radii, give $\abs a/\abs t<1/(2T)$ for $T=18.08$ and hence $d\le327$ \cite{LMFDBDatabase}. Ng's
verification of RH for this Artin factor up to height $972.2$ \cite[\S\S4.3.4--4.3.5]{Ng2000Limiting} gives
$d\le945{,}173$. These computations were not repeated.

\paragraph{Cubic Jensen tests near the Gaussian.}
The curvature inequality $(\log K(\sqrt x))''<0$ used in \cref{hec:thm:turan} does not imply the cubic
inequalities, contrary to what Theorem~1 of \cite{DimitrovLucas2011Higher} would give for admissible
kernels. The reason is visible at first order around the Gaussian, and the first-order analysis has exact
answers. Let $V\colon[0,\infty)\to\R$ be continuous, bounded below and of at most polynomial growth, and for
$\eps>0$ let
\[
 K_\eps(u)=\exp\bigl[-u^2-\eps V(u^2)\bigr].
\]
Let $X_n$ have the gamma distribution of shape $n+1/2$ and rate one, and put
\[
 m(n)=\hecE V(X_n),\qquad\Delta m(n)=m(n+1)-m(n).
\]
Call a shift $n$ \emph{persistent} for $V$ if $J^{3,n}$, formed with the kernel $K_\eps$, is real-rooted
along some sequence $\eps\downarrow0$; each shift may use its own sequence.

\begin{lemma}[First-order discriminant]\label{hec:lem:disc}
With $c_n=4^n\gamma_n/\sqrt\pi$, one has $c_n=\hecE e^{-\eps V(X_n)}=1-\eps m(n)+O_n(\eps^2)$. The cubic
$P(W)=\sum_{j=0}^3\binom3jc_{n+j}(W-1)^j$, which differs from $J^{3,n}$ by positive rescalings of the
variable and of the value, has discriminant
\[
 -27\eps^2\bigl(\Delta^3m(n)\bigr)^2+O_n(\eps^3)\quad\text{if }\Delta^3m(n)\ne0,\qquad
 108\eps^3\bigl(\Delta^2m(n+1)\bigr)^3+O_n(\eps^4)\quad\text{if }\Delta^3m(n)=0 .
\]
In particular a persistent shift has $\Delta^3m(n)=0$.
\end{lemma}

\begin{proof}
The substitution $x=u^2$ and the duplication formula for the gamma function give the formula for $c_n$.
The expansion follows by integrating the second-order Taylor remainder of the exponential, which is
bounded by a constant times $\eps^2V^2$ for $0<\eps<1$ because $V$ is bounded below. The coefficients of
$1,W,W^2,W^3$ in $P$ are, respectively,
\[
 \eps\Delta^3m(n)+O_n(\eps^2),\qquad-3\eps\Delta^2m(n+1)+O_n(\eps^2),\qquad O_n(\eps),\qquad1+O_n(\eps),
\]
and the two expansions of the discriminant follow.
\end{proof}

\begin{theorem}[Polynomial perturbations: the number of tests; \stU]\label{hec:thm:finitejensen}
Fix $r\ge3$ and let $V$ be a polynomial of degree at most $r$, bounded below on $[0,\infty)$. If $r-2$
distinct shifts are persistent, then $\deg V\le2$, and in that case every $J^{d,n}$, $d\ge1$, $n\ge0$, has
only negative real roots for every $\eps>0$. The number $r-2$ is sharp: for every set $S$ of $r-3$
nonnegative integers there is a strictly convex increasing polynomial of degree $r$ whose persistent shifts
are exactly the elements of $S$, and at each fixed shift the predicted outcome holds for all sufficiently
small $\eps>0$.
\end{theorem}

\begin{proof}
The map $V\mapsto m$ sends $x^k$ to the monic polynomial $(n+1/2)(n+3/2)\cdots(n+k-1/2)$ in $n$, so it is
a degree-preserving linear isomorphism. If $r-2$ shifts are persistent, the polynomial $\Delta^3m$, of
degree at most $r-3$, vanishes at $r-2$ points by \cref{hec:lem:disc}, hence identically, and
$\deg V\le2$.

Conversely let $V(x)=ax^2+bx+c$. For $a>0$, P\'olya's theorem on quartic Fourier integrals with an
arbitrary quadratic coefficient gives only real zeros for $\int e^{itu}K_\eps(u)\,\dd u$
\cite[(1.3)]{deBruijn1950Roots}. Then $G(w)=\int\cosh(\sqrt w\,u)K_\eps(u)\,\dd u$ has only negative real
zeros, and the bound $\abs{\sqrt w}\abs u\le(\eps a/2)u^4+C\abs w^{2/3}$ shows that its order is at most
$2/3$, so its genus-zero product places it in the Laguerre--P\'olya class. For $a=0$, boundedness below
gives $b\ge0$ and $G$ is a positive multiple of an exponential with positive slope, again in the class.
Every derivative of $G$ is in the class, so Jensen's theorem gives real roots for every $J^{d,n}$, and
positivity of the coefficients puts them on the negative axis.

For sharpness take $m$ with $\Delta^3m(n)=\prod_{s\in S}(n-s)$, which is possible by triangularity of
finite differences. Its leading coefficient, and that of the corresponding $V$, is positive. Adding $Ax^2$
to $V$ adds $2A$ to $\Delta^2m(n+1)$ and preserves $\Delta^3m$. Choose $A$ so large that
$\Delta^2m(s+1)>0$ for each $s\in S$ and $V''>0$ on $[0,\infty)$, and add a linear term to make $V'(0)>0$.
The two expansions in \cref{hec:lem:disc} give all claims; the threshold for $\eps$ may depend on the
shift.
\end{proof}

\begin{theorem}[The exact stride; \stU]\label{hec:thm:stride}
Let $V$ be continuous, bounded below and of at most polynomial growth. If all shifts in $r+2\N_0$ are
persistent for some integer $r\ge0$, then $V$ is a polynomial of degree at most two. For every stride
$k\ge3$ and every residue $r\bmod k$ there is a smooth, strictly convex, increasing, nonpolynomial $V$
whose persistent shifts are exactly that residue class.
\end{theorem}

Thus one parity class of shifts, even with a finite initial segment missing, detects every nonquadratic
perturbation in this class; in particular, real-rootedness of all cubic Jensen polynomials along a sequence
$\eps\downarrow0$ forces $V$ to be a polynomial of degree at most two.

\begin{proof}
By \cref{hec:lem:disc}, persistence gives $\Delta^3m(n)=0$ on the progression. Define a distribution
supported on $[0,\infty)$ by
\[
 T(\phi)=-\int_0^\infty V(x)\frac{\dd^3}{\dd x^3}\bigl(x^{r+5/2}e^{-x}\phi(x)\bigr)\,\dd x .
\]
The least exponent at the origin is $r-1/2>-1$, and by polynomial growth and exponential decay $T$ can be
tested on polynomials and on the exponentials $e^{zx}$ with $\Re z<1$. For $j\ge0$ and $p=r+j+1/2$,
differentiating the integrand gives
\[
 T(x^j)=\Gamma(p+3)\Delta^3m(r+j),
\]
because the four terms have integrals proportional to $m(r+j+3),-3m(r+j+2),3m(r+j+1),-m(r+j)$ with the same
factor. Put $F(z)=T(e^{zx})$, analytic in $\Re z<1$. If the progression of stride two is persistent, all
even derivatives of $F$ at zero vanish, so $F(z)+F(-z)=0$ for $\abs{\Re z}<1$. On the imaginary axis this
says that $T$ plus its reflection has zero Fourier transform; both are tempered, so
$T+T^{\mathrm{refl}}=0$. Their supports meet only at the origin, so $T$ vanishes on $(0,\infty)$, where
$T=x^{r+5/2}e^{-x}V'''$ in the sense of distributions. The multiplier is smooth and positive, so $V'''=0$
and $V$ is a polynomial of degree at most two. \Cref{hec:thm:finitejensen} gives the converse.

For sharpness set $t=\tan(\pi/k)>0$ and $h=(1-it)^{-1}$, so that $\arg h=\pi/k$, and take
\[
 V(x)=Ax^2+Bx+\cos(tx+\theta).
\]
The gamma expectation of the oscillatory part is $\Re(e^{i\theta}h^{n+1/2})$, and its third difference is
$\Re(e^{i\theta}h^{n+1/2}(h-1)^3)$. Choose $\theta=\pi/2-(r+1/2)\pi/k-3\arg(h-1)$. Since $h\ne1$, the third
difference vanishes exactly for $n\equiv r\pmod k$; the quadratic and linear terms do not change it. The
oscillatory contribution to the second difference has absolute value at most $\abs{h-1}^2$, since
$\abs h<1$. Choose $2A>\max(t^2,\abs{h-1}^2)$ and $B>t$. Then $V'>0$, $V''>0$ and $\Delta^2m(n+1)>0$
everywhere. On the chosen progression the leading term of the discriminant is positive of order $\eps^3$,
and off it the leading term is negative of order $\eps^2$.
\end{proof}

\begin{example}[Explicit cubic failures]\label{hec:ex:sextic}
(a) For $K_\delta(u)=\exp(-u^2-\delta u^6)$, that is $V(x)=x^3$, the third difference of
$m(n)=(n+1/2)(n+3/2)(n+5/2)$ is $6$, and the discriminant of the rescaled cubic at shift $0$ is
$-972\delta^2+O(\delta^3)$. For $\delta=10^{-8}$ the Taylor inequalities
$1-\delta(p)_3\le c_n\le1-\delta(p)_3+\tfrac12\delta^2(p)_6$, with $p=n+1/2$ and $(p)_j$ the rising
factorial, certify
\[
 -975<\disc(P)/\delta^2<-969
\]
(\artifact{grh_transfers_density_gaps_and_zero_tests/checks/check_cubic.py}), so $J^{3,0}$ has a nonreal
pair. This kernel is positive, even, entire, strictly decreasing for $u>0$, rapidly decreasing with all
derivatives, and satisfies $(\log K_\delta(\sqrt x))''=-6\delta x<0$; it is a counterexample to the
general cubic criterion. In the argument of \cite{DimitrovLucas2011Higher} the difficulty is in the
mean-value step preceding their equation~(14): the function
$A(y,z)=(z^2-y^2)K(y)K(z)[B(y)-B(z)]$, with $B(y)=K'(y)/(yK(y))$, vanishes to second order as $z\to y$, so
$\partial_yA/A$ has a simple pole at $z=y$ and the step cannot be applied on an interval with
$z_1<y<z_2$.
{\emergencystretch=14em\par}

(b) The obstruction persists under strict fourth-derivative curvature. For
$\widehat K_\delta(u)=\exp[-u^2-\delta(u^4+u^6)]$ the logarithm has fourth derivative
$-24\delta-360\delta u^2<0$, as for the eta kernels of \cref{hec:thm:turan}, and $V(x)=x^2+x^3$ again has
$\Delta^3m=6$. A single sufficiently small $\delta$ violates the cubic inequalities at any prescribed
finite set of shifts.

(c) For $V(x)=x^4-14x^3+100x^2+x$ one has $V''(x)=12(x-7/2)^2+53$, $\Delta^3m(n)=24n$ and
$\Delta^2m(1)=95$, so shift $0$ is persistent and every other fixed shift fails for small $\eps$. At
$\eps=10^{-12}$ the inequalities $1-y\le e^{-y}\le1-y+y^2/2$ give
\[
 92{,}590{,}000<\disc(P_0)/\eps^3<92{,}600{,}000,\qquad-15{,}553<\disc(P_1)/\eps^2<-15{,}551
\]
(\artifact{grh_dihedral_47_71/check_finite.py}).
{\emergencystretch=14em\par}
\end{example}

The first-order analysis leaves open whether a quartic term can restore the low-degree tests while a
higher degree still detects nonreal zeros. This is decided by a second scaling. Fix $t>0$ and $c\ge0$,
and consider
\[
 K_{t,c}(u)=\exp(-u^2-t^2u^4-ct^3u^6),\qquad
 b_n=4^n\gamma_n/\sqrt\pi,\qquad J_{d,n}(Z)=\sum_{k=0}^d\binom dkb_{n+k}Z^k .
\]

\begin{theorem}[Sextic perturbations; \stU]\label{hec:thm:sexticlimit}
Let $D=\dd/\dd y$ and $Q_{d,c}(y)=\exp(-D^2+cD^3)\,y^d$.
\begin{enumerate}[label=\textup{(\roman*)}]
\item For fixed $n,d,c$, $t^{-d}J_{d,n}(-1+ty)\to Q_{d,c}(y)$ coefficientwise as $t\downarrow0$.
\item Let $c_*=2\sqrt2/3$. For every fixed shift, $J_{3,n}$ has three distinct real roots for all
sufficiently small $t>0$ if $0\le c\le c_*$, and a nonreal pair for all sufficiently small $t>0$ if
$c>c_*$.
\item If $d\ge3$ and $Q_{d,c}$ is real-rooted, then $c^2\le8/(9(d-2))$. Consequently, for $c>0$, any
$d>2+8/(9c^2)$ and any fixed $n$, the polynomial $J_{d,n}$ has a nonreal pair for all sufficiently small
$t>0$. Conversely, for any finite set of degrees and shifts one can first fix a small $c>0$ and then take
$t>0$ so small that every prescribed $J_{d,n}$ has distinct real roots.
\end{enumerate}
All this holds although $(\log K_{t,c}(\sqrt x))''=-2t^2-6ct^3x<0$ everywhere.
\end{theorem}

\begin{proof}
(i) With $p=n+1/2$ and $X$ gamma distributed of shape $p$, $b_n=\hecE H(tX)$ where
$H(s)=\exp(-s^2-cs^3)$. If $h_j$ is the coefficient of $s^j$ in $H$, Taylor's formula gives
$b_n=\sum_{j=0}^dh_jt^j(p)_j+O_{n,d,c}(t^{d+1})$, since every derivative of $H$ is bounded on the positive
half-line and $X$ has all moments. At index $n+k$, $(p+k)_j$ is monic of degree $j$ in $k$. Acting on
$(1+Z)^d$ with $(Z\partial_Z)^j$ inserts the multiplier $k^j$ in the binomial sum, and the top derivative
in this operator is $Z^j\partial_Z^j$, which at $Z=-1+ty$ yields
$(-1)^jd(d-1)\cdots(d-j+1)\,t^{d-j}y^{d-j}+O(t^{d-j+1})$; all lower derivatives carry at least one more
power of $t$. Summing gives $H(-D)y^d$. The polynomials $Q_{d,c}$ are classical three-variable Hermite
polynomials; a broader classification of Fourier transforms with only real zeros under Gaussian damping
is in \cite{Newman1976Fourier}.

(ii) $Q_{3,c}=y^3-6y+6c$ has discriminant $864-972c^2$, which settles the case $c\ne c_*$ by continuity of
the roots. At $c_*$ the limit has a double root at $y=\sqrt2$. The first five Taylor coefficients
$(1,0,-1,-c,1/2)$ of $H$ give
\[
 t^{-3}J_{3,n}(-1+ty)=Q_{3,c}(y)+6t(p+3)(y^2-3cy-2)+O_{n,c}(t^2),
\]
and the correction at the double root is $-24t(p+3)$. The local minimum is therefore negative, the values
at a fixed small distance on either side are positive, and the simple third root persists. More precisely,
for $c=c_*+at$ the local minimum is $t(6a-24(p+3))+O(t^2)$, so the value $a=4(p+3)$ separates the two
outcomes to this order.

(iii) Let $r_i$ be the roots of $Q_{d,c}$ and $s_k=\sum r_i^k$. Newton's identities with the first four
coefficients give
\[
 s_1=0,\qquad s_2=2d(d-1),\qquad s_3=-3cd(d-1)(d-2),\qquad s_4=4d(d-1)(2d-3).
\]
For real $r_i$ the Cauchy--Schwarz inequality for $r_i$ and $r_i^2-s_2/d$ gives
$s_3^2\le s_2(s_4-s_2^2/d)$, which is the stated bound. The remaining assertions follow from (i), since
$Q_{d,0}$ is a rescaled Hermite polynomial with simple real zeros.
\end{proof}

The bound in (iii) is sufficient for failure; its optimal value for each $c$ is not determined. At $n=0$,
$t=1/10000$, $c=4/5$ the cubic passes and the quartic fails: bounding each $b_k$ by the Taylor polynomials
of degrees nine and ten of $e^{-V}$, $V=t^2X^2+ct^3X^3$, averaged over the gamma law, the rescaled cubic
has alternating signs at $-4,0,7/5,3$, and for the quartic the determinant of $(s_{i+j})_{0\le i,j\le2}$
lies strictly between $-11000$ and $-1000$, whereas real roots would make this matrix a Gram matrix
(\artifact{mcnew_density_and_grh_transfers/checks/verify_jensen.py}).
{\emergencystretch=14em\par}

These results describe the reach of arguments that use only convexity properties of the kernel. Real
roots in every degree for the conductor-$23$ kernel would require an argument uniform in the degree that
uses arithmetic information.

\subsection{Class-group Hecke \texorpdfstring{$L$}{L}-functions}\label{hec:ss:classgroup}
{\emergencystretch=\textwidth\sources{grh_method_transfers_sharp_torsion}}

In the source proof the residues of the prime-error transform are differences of zero multiplicities, and
an auxiliary twist is chosen to remove the competing term. For Hecke $L$-functions of class characters one
can instead choose weights on ideal classes so that the unwanted unknown factors disappear from the
logarithmic derivative and the remaining residues have a common sign.

Let $K$ be an imaginary quadratic field, $C$ its ideal class group, $h=\abs C$, and $B\le C$ a subgroup.
Let $B^\perp$ be the group of characters of $C$ trivial on $B$, and suppose that the functions
$L_K(s,\eta)$, $\eta\in B^\perp$, have no zeros in $\Re s>1/2$; these are the \emph{known factors}. For
$B=C^2$ every such character has order at most two, Kronecker's genus factorization writes its
$L$-function as a product of two Dirichlet $L$-functions (for the trivial character,
$\zeta_K=\zeta L(\cdot,\chi_K)$), and \cref{src:thm:dgrh} supplies the hypothesis. For background on
class group $L$-functions see \cite{DukeFriedlanderIwaniec1995Class}.

For a real, inversion-invariant weight $w\colon C\to[0,1]$ put
\[
 a_\chi=\frac1h\sum_{c\in C}w(c)\overline{\chi(c)},\qquad
 \Theta_w(x)=\sum_{\substack{p\le x\\p\ \mathrm{split\ in}\ K}}w([\mathfrak p])\log p,\qquad
 \delta=\frac{a_1}2,
\]
where $a_1$ is the coefficient of the trivial character. Either prime ideal above $p$ gives the same
weight, each rational prime is counted once, and the $a_\chi$ are real. Let $S$ be a nonempty
inversion-invariant subset of $\widehat C\setminus B^\perp$, and require
\begin{equation}\label{hec:eq:weightcond}
 a_\chi<0\quad(\chi\in S),\qquad a_\chi=0\quad(\chi\notin S\cup B^\perp).
\end{equation}
Let $P_w(s)=\sum_{p\ \mathrm{split}}w([\mathfrak p])(\log p)p^{-s}$,
$c_w=\lim_{s\to1}(P_w(s)-\delta/(s-1))-\delta$, and
\[
 R_w(x)=\sum_{\substack{p\le x\\p\ \mathrm{split\ in}\ K}}w([\mathfrak p])\log p\,
 \Bigl(\frac1p-\frac1x\Bigr)-\delta\log x-c_w .
\]

\begin{theorem}[Class-weight criterion; \stR]\label{hec:thm:classweight}
With the field, the subgroup and the weight fixed as above, the following are equivalent.
\begin{enumerate}[label=\textup{(\arabic*)}]
\item Each $L_K(s,\chi)$ with $\chi\in S$ satisfies RH.
\item For each $0<\eps<1/2$ there is $M_\eps$ with $\Theta_w(x)-\delta x\le M_\eps x^{1/2+\eps}$ for all
sufficiently large $x$.
\item For each $0<\eps<1/2$ there is $M_\eps$ with $R_w(x)\ge-M_\eps x^{-1/2+\eps}$ for all sufficiently
large $x$.
\end{enumerate}
\end{theorem}

\begin{proof}
Comparing Euler factors for $\Re s>1$ gives
\[
 P_w(s)=-\frac12\sum_\chi a_\chi\frac{L_K'}{L_K}(s,\chi)+H(s)
\]
with $H$ holomorphic in $\Re s>1/2$: the two primes of degree one above a split prime contribute equally,
inert primes have norm $p^2$, the powers of exponent at least two converge absolutely and locally
uniformly there, and the ramified primes contribute finitely many locally holomorphic series. The pole at
$1$ has residue $\delta$. At a zero $\rho$ with $1/2<\Re\rho<1$ of the tested functions the residue of
$P_w$ is
\[
 d_\rho=\frac12\sum_{\chi\in S}(-a_\chi)\ord_\rho L_K(s,\chi)>0 ;
\]
an untested function vanishing at the same point does not enter, and nothing is assumed about simplicity
or independence of zeros. With $E(x)=\Theta_w(x)-\delta x$, termwise integration in $\Re s>1$ gives
\begin{align*}
 \int_1^\infty E(x)x^{-s-1}\,\dd x&=\frac{P_w(s)}s-\frac\delta{s-1},\\
 \int_1^\infty xR_w(x)x^{-s-1}\,\dd x&=\frac{P_w(s)}{s(s-1)}-\frac\delta{(s-1)^2}-\frac{c_w}{s-1},
\end{align*}
both regular at $1$. Their residues at a real zero $\rho\in(1/2,1)$ are $d_\rho/\rho>0$ and
$d_\rho/(\rho(\rho-1))<0$. Let $a=1/2+\eps$. Under (2), \cref{hec:lem:landau} applies to $g=E$; under (3)
it applies to $g=-xR_w$. In both cases the transform is holomorphic in $\Re s>a$, so no tested function
has a zero there. Vary $a$ and use the functional equations, $S$ being closed under inversion.

Conversely, under (1) and the hypothesis on the known factors, the explicit formula gives
$\sum_{\Nm\mathfrak p^k\le x}\chi(\mathfrak p)^k\log\Nm\mathfrak p
=\hecind_{\chi=1}x+O_{K,\chi}(\sqrt x\log^2(2x))$ for each $\chi$ with $a_\chi\ne0$. Removing prime powers,
inert primes and ramified primes costs $O_K(\sqrt x\log^2(2x))$, and Fourier inversion gives the same
two-sided error for $\Theta_w$, which is (2). Finally $\sum_{p\le x}w\log p\,(1/p-1/x)
=\int_1^x\Theta_w(t)t^{-2}\,\dd t$; the two-sided error makes $\int_1^\infty E(t)t^{-2}\,\dd t$ converge with
tail $O_{K,w}(x^{-1/2}\log^2(2x))$, and evaluating the first transform at $s=1$ identifies the integral
with $c_w$. This gives (3).
\end{proof}

For $w=\hecind_{B\setminus\set1}$, orthogonality gives $a_\eta=(\abs B-1)/h$ for $\eta\in B^\perp$ and
$a_\chi=-1/h$ otherwise, so the test detects every remaining factor, with prime density
$(\abs B-1)/(2h)$. In a ring class field the event is splitting completely in the fixed field of $B$ and
not in the full field. This aggregate case is an instance of the normal-subgroup criterion of
\cref{sec:nab}. Selecting a single pair of factors requires a different weight.

\paragraph{Isolating one dihedral factor.}
Fix a character $\chi\notin B^\perp$ of order $n>2$. Since $L_K(s,\chi)=L_K(s,\chi^{-1})$, we seek a
weight with $a_\chi=a_{\chi^{-1}}=-t<0$ and coefficient zero on every other unknown character; by
\cref{hec:thm:classweight} it tests this Hecke $L$-function alone, even if other untested functions have
zeros off the line.

\begin{theorem}[Sharp coefficient, support and coefficient field; \stU]\label{hec:thm:isolate}
Put $v(c)=\Re\chi(c)$, and for each coset $Q$ of $B$ let $M_Q=\max_{c\in Q}v(c)$,
$m_Q=\min_{c\in Q}v(c)$ and $D=\max_Q(M_Q-m_Q)$.
\begin{enumerate}[label=\textup{(\roman*)}]
\item $D>0$, the largest possible $t$ is $1/(2D)$, and it is attained by
\[
 w_\chi(c)=\frac{M_{cB}-v(c)}D .
\]
This weight also has the smallest support, by inclusion, among all admissible weights.
\item For $B=C^2$,
\[
 t_{\max}=\begin{cases}
 [2(1+\cos(\pi/n))]^{-1},&n\ \text{odd},\\
 1/4,&4\mid n,\\
 [2(1+\cos(2\pi/n))]^{-1},&n\equiv2\pmod4 .
 \end{cases}
\]
\item The smallest real number field containing all values of a weight that isolates the pair
$\set{\chi,\chi^{-1}}$ is $F_n=\Q(\zeta_n+\zeta_n^{-1})$.
\item For $B=C^2$, a rational-valued weight isolating the pair exists exactly for $n\in\set{3,4,6}$, and
in these cases a $\set{0,1}$-valued weight suffices.
\end{enumerate}
\end{theorem}

\begin{proof}
(i) Fourier inversion gives $w(c)=u(cB)-2tv(c)$, where $u$ is a real inversion-invariant function on
$C/B$. The bounds $0\le w\le1$ hold exactly when $2tM_Q\le u(Q)\le1+2tm_Q$ for every coset $Q$. Since
$\chi|_B$ is nontrivial, $v$ is not constant on $B$, so $D>0$ and $t\le1/(2D)$. Taking $u(Q)=2tM_Q$ makes
the weight vanish at every coset maximum, while all other points have positive weight for every feasible
$u$ and every $t>0$.

(ii) For odd $n$ the image of $B$ under $\chi$ consists of all $n$th roots of unity, with extreme real
parts $1$ and $-\cos(\pi/n)$. For even $n$ each coset maps onto the even or the odd powers of a primitive
$n$th root. The even powers include $\pm1$ when $4\mid n$. For $n\equiv2\pmod4$ they form a regular polygon
with an odd number of vertices, and the odd powers are its negative.

(iii) The weight $w_\chi$ has values in $F_n$. Conversely, let $E$ be a number field containing all values
of an isolating weight. Every automorphism of $\overline\Q$ fixing $E$ induces a power map on $\mu_n$ and
permutes the nonzero Fourier coefficients belonging to primitive powers of $\chi$. These powers are all
nontrivial on $B$, because raising a root of unity to a power prime to $n$ preserves its order. Hence the
image of $\chi$ is $\chi$ or $\chi^{-1}$, every such automorphism fixes $\zeta_n+\zeta_n^{-1}$, and
$F_n\subseteq E$.

(iv) An automorphism taking $\zeta_n$ to $\zeta_n^j$ sends $a_\chi$ to $a_{\chi^j}$, so a rational weight
has nonzero coefficients on the whole orbit of $\chi$, all of whose members have order $n>2$ and are
unknown. An isolated pair forces $\varphi(n)\le2$. For sufficiency write $\chi(c)=\exp(2\pi ij/n)$ and take
the indicators of the exponents $\set{1,2}$ for $n=3$, $\set2$ for $n=4$, and $\set{2,3,4}$ for $n=6$;
these are the weights $w_\chi$.
\end{proof}

The classification concerns exact vanishing of the extraneous Fourier coefficients; it says nothing about
coincidences of zeros.

\begin{example}[A pair that requires irrational weights]\label{hec:ex:d95}
For $C=\Z/8\Z$ and $\chi(j)=\exp(2\pi ij/8)$ the optimal weight is, in the order $j=0,\dots,7$,
\[
 (0,0,1/2,\sqrt2/2,1,\sqrt2/2,1/2,0).
\]
Its coefficient on each of $\chi,\chi^{-1}$ is $-1/4$ and its rational-prime mean is $(2+\sqrt2)/16$. The
primitive character $\chi^3$ gives a second pair, which does not contribute; replacing $j$ by $3j$ gives
its own test. A rational weight could detect both pairs together but cannot separate them.

A realization by a maximal order is $K=\Q(\sqrt{-95})$, whose eight reduced primitive forms are
\[
 [1,1,24],\quad[2,\pm1,12],\quad[3,\pm1,8],\quad[4,\pm1,6],\quad[5,5,6].
\]
Let $g=[2,1,12]$. The root $b=1$ of $b^2\equiv-95$ modulo $4\cdot2^k$ for $k=1,2,3$ gives $g^2=[4,1,6]$ and
$g^3=[8,1,3]=[3,-1,8]$ by composition. The square $[4,1,6]$ differs from its inverse $[4,-1,6]$, so $g$ has
order eight, and the only nonprincipal ambiguous class $[5,5,6]$ is $g^4$. The test therefore uses the
three representation conditions
\begin{center}
\begin{tabular}{c|ccc}
$Q(x,y)$ & $4x^2+xy+6y^2$ & $3x^2+xy+8y^2$ & $5x^2+5xy+6y^2$\\\hline
$w$ & $1/2$ & $\sqrt2/2$ & $1$
\end{tabular}
\end{center}
on primes $p\nmid95$; for prime inputs the three sets are disjoint. Replacing the middle form by
$2x^2+xy+12y^2$ gives the test for the other faithful pair
(\artifact{grh_method_transfers_sharp_torsion/checks/verify_d95.py}).
{\emergencystretch=14em\par}
\end{example}

\paragraph{The Gaussian ring class tower.}
Let $K=\Q(i)$, let $K_n$ be its ring class field of conductor $2^n$ for $n\ge1$, and for $n\ge2$ let
\[
 Q_n(X,Y)=(2X+Y)^2+2^{2n-2}Y^2,\qquad
 \vartheta_n(x)=\sum_{\substack{p\le x\\ p=Q_n(a,b)}}\log p,
\]
each represented odd prime being counted once.

\begin{theorem}[One quadratic form for each layer; \stR]\label{hec:thm:tower}
Let $n\ge2$ and assume RH for $\zeta_{K_{n-1}}$. Then RH for $\zeta_{K_n}$ is equivalent to
\[
 \vartheta_n(x)\le2^{-n}x+O_{n,\eps}(x^{1/2+\eps})\qquad\text{for every }\eps>0 .
\]
Under \cref{src:thm:dgrh}, these bounds for all $n\ge3$ together are equivalent to RH for every
$\zeta_{K_n}$.
\end{theorem}

\begin{proof}
The ring class number formula, with the factor two from the units, gives $[K_n:K]=2^{n-1}$, and an odd
prime splits completely in $K_n$ exactly when $p=a^2+2^{2n}b^2$ \cite[\S\S7, 9]{Cox2013Primes}. The fields
are nested. For $p\equiv1\pmod4$ write $p=a^2+c^2$ with $a$ odd and $c$ even; the pair of absolute values
is unique. Complete splitting in $K_{n-1}$ together with failure of complete splitting in $K_n$ means
$v_2(c)=n-1$, that is, $c=2^{n-1}Y$ with $Y$ odd and $a=2X+Y$; conversely an odd prime of the form $Q_n$
forces $Y$ odd and this valuation. So $\vartheta_n$ counts exactly these primes. The quadratic extension
$K_2=\Q(\zeta_8)$ of $K$ is abelian over $\Q$, so \cref{src:thm:dgrh} gives RH for $\zeta_{K_2}$ and starts
the induction.

For a Galois field $M/\Q$, the series $P_M(s)=\sum_{p\ \text{split completely}}(\log p)p^{-s}$ equals
$-[M:\Q]^{-1}\zeta_M'(s)/\zeta_M(s)$ up to a function holomorphic in $\Re s>1/2$, because the unramified
primes of degree one occur in full packets of $[M:\Q]$ and all other terms have exponent at least two. The
tested series is $P_{K_{n-1}}-P_{K_n}$ up to finitely many primes; its pole at $1$ has residue
$2^{1-n}-2^{-n}=2^{-n}$, and under RH for the lower field a zero of $\zeta_{K_n}$ at $\beta+i\gamma$ with
$1/2<\beta<1$ contributes a positive residue, its multiplicity divided by $2^n$. \Cref{hec:lem:landau},
applied to $\vartheta_n(x)-2^{-n}x$, gives zero-freeness from the bound. Conversely RH for both fields and
the prime ideal theorem give the two-sided error $O_n(\sqrt x\log^2(2x))$. No holomorphy of individual
Artin factors is used.
\end{proof}

The first case is $Q_3=4X^2+4XY+17Y^2$, of density $1/8$; here $K_3=\Q(i,\sqrt[4]2)$ and the test counts
the primes $p\equiv1\pmod8$ for which $2$ is a quartic nonresidue. Detecting $v_2(c)=2$ in $p=a^2+c^2$
uses splitting in a nonabelian extension of $\Q$: the class group at conductor eight is cyclic of order
four and complex conjugation inverts it, so the Galois group over $\Q$ is dihedral of order eight, and the
test is not a restriction to Dirichlet residue classes.

\subsection{Hecke characters of the Gaussian field}\label{hec:ss:gaussian}
{\emergencystretch=\textwidth\sources{grh-methods-polynomial-density}}

Let $K=\Q(i)$. Every nonzero ideal has a generator $a$, and $\eta((a))=(a/\abs a)^4$ does not depend on its
choice. Write $L_k(s)=L_K(s,\eta^k)$ and $L_0=\zeta_K$. For $k\ne0$ these are entire Hecke $L$-functions of
nonzero infinity type, with $L_{-k}=L_k$; their zeros govern the distribution of the angles of Gaussian
primes \cite{RudnickWaxman2017Angles}. For a real even trigonometric polynomial
$P(z)=h_0+\sum_{k=1}^Nh_k(z^k+z^{-k})$ on the unit circle,
\begin{equation}\label{hec:eq:DP}
 D_P(s)=\sum_{\mathfrak p,\,m\ge1}\frac{\log\Nm\mathfrak p\;P(\eta(\mathfrak p)^m)}{(\Nm\mathfrak p)^{ms}}
 =-h_0\frac{\zeta_K'}{\zeta_K}(s)-2\sum_{k=1}^Nh_k\frac{L_k'}{L_k}(s)\qquad(\Re s>1),
\end{equation}
the sum being over all prime ideals. A nonnegative $P$ gives nonnegative local weights, nonnegative $h_k$
make the orders of zeros add with the same sign, and the mean $h_0$ is the residue of $D_P$ at $1$. One is
led to ask for the smallest main term compatible with prescribed sensitivity to several $L_k$.

\begin{theorem}[Optimal weights for a block of modes; \stU]\label{hec:thm:belov}
Fix $0<\alpha<1$ and integers $1\le r\le N$. With rising factorials $(a)_j=a(a+1)\cdots(a+j-1)$ put
\[
 c_j=\frac{(\alpha)_j}{j!},\qquad b_j=\frac{(2\alpha)_j}{j!},\qquad
 d=\min\set{r,\lfloor N/2\rfloor+1}.
\]
Among nonnegative $P$ of degree at most $N$ with $h_{N-j}\ge b_j$ for $0\le j<r$, the minimum of $h_0$ is
\[
 M_{N,r}(\alpha)=2\sum_{j=0}^{d-1}c_j^2-\hecind_{\set{N=2d-2}}\,c_{d-1}^2 .
\]
It is attained with every $h_k\ge0$, by $P(z)=\abs{\sum_{i=0}^Nv_iz^i}^2$ with $v_i=c_{\min(i,N-i)}$ if
$\min(i,N-i)<d$ and $v_i=0$ otherwise.
\end{theorem}

\begin{proof}
The $c_i$ decrease and $(1-z)^{-\alpha}=\sum c_iz^i$, so $\sum_{i=0}^jc_ic_{j-i}=b_j$. By the symmetry of
$v$, $h_{N-j}=\sum_{i=0}^jv_iv_{j-i}$. For $j<d$ this equals $b_j$. If $d\le j<r$, the support covers every
index and $v_i\ge c_i$ throughout, so the remaining constraints hold. Parseval's identity gives the stated
mean, and every autocorrelation $h_k$ is nonnegative.

For the lower bound use the Fej\'er--Riesz factorization $P(z)=\abs{\sum_{i=0}^Nu_iz^i}^2$. Put
$A_i=\abs{u_i}$ and $B_i=\abs{u_{N-i}}$ for $0\le i<d$, and $w_i=1$ except that $w_{d-1}=1/2$ when
$N=2d-2$; then $\sum_iw_i(A_i^2+B_i^2)\le h_0$, the centre index being counted once. Let $e_k$ be the
coefficient of $z^k$ in $(1-z)^\alpha$, set $y_i=w_ic_i$, and define
\[
 \lambda_j=\sum_{k=0}^{d-1-j}e_ky_{j+k},\qquad
 H_{ij}=\begin{cases}\lambda_{i+j},&i+j<d,\\0,&i+j\ge d.\end{cases}
\]
Here $e_0=1$, $e_k<0$ for $k\ge1$, and $\sum_{k=0}^me_k>0$, the coefficient of $z^m$ in
$(1-z)^{\alpha-1}$. Since $y_i$ is positive and nonincreasing, $\lambda_j>0$. The convolution of $c$ with
$e$ is $(1,0,\dots)$, which gives $(Hc)_i=\sum_{j<d-i}c_j\lambda_{i+j}=w_ic_i$. For all $i,j$ the
arithmetic--geometric mean inequality gives $2A_iB_j\le A_i^2c_j/c_i+B_j^2c_i/c_j$. Multiplying by
$H_{ij}$, summing, and using the symmetry of $H$ and $Hc=(w_ic_i)_i$,
\[
 2\sum_{j<d}\lambda_jb_j\le2\sum_{j<d}\lambda_jh_{N-j}\le2\sum_{i,j<d}H_{ij}A_iB_j
 \le\sum_{i<d}w_i(A_i^2+B_i^2)\le h_0 .
\]
The first sum, without the factor two, is $\sum_{i,j<d}H_{ij}c_ic_j=\sum_{i<d}w_ic_i^2$, which is half of
$M_{N,r}(\alpha)$.
\end{proof}

For $\alpha=1/2$ the required gains are $b_j=1$, and with $r=N$ the theorem is Belov's, after
multiplication of his cosine polynomial by two; Belov also proved uniqueness in that case
\cite{Belov1994Extremal}. For $N=r=4$ and $\alpha=1/3$ the factor is $1+z/3+2z^2/9+z^3/3+z^4$ and the
optimal mean is $184/81$. The constraints at the lower frequencies beyond $\lfloor N/2\rfloor+1$ never
increase the optimum. Since $c_j\sim j^{\alpha-1}/\Gamma(\alpha)$, the full block satisfies, as
$N\to\infty$,
\begin{equation}\label{hec:eq:belovasym}
 M_{N,N}(\alpha)=
 \begin{cases}
 2\sum_{j\ge0}c_j^2+o(1),&0<\alpha<1/2,\\[2pt]
 (2/\pi)\log N+O(1),&\alpha=1/2,\\[2pt]
 \dfrac{2(N/2)^{2\alpha-1}}{(2\alpha-1)\Gamma(\alpha)^2}\,(1+o(1)),&1/2<\alpha<1 .
 \end{cases}
\end{equation}
In the bounded regime the gains decrease, but each of the $N$ gains is positive at every fixed $N$.

Let $g_P=\lim_{s\to1}(D_P(s)-h_0/(s-1))-h_0$ and
\[
 U_P(x)=\sum_{(\Nm\mathfrak p)^m\le x}\log\Nm\mathfrak p\;P(\eta(\mathfrak p)^m)
 \Bigl(\frac1{(\Nm\mathfrak p)^m}-\frac1x\Bigr)-h_0\log x-g_P .
\]

\begin{theorem}[Angular criterion; \stR]\label{hec:thm:angular}
Let $P$ be the extremal polynomial of \cref{hec:thm:belov}. If for each $0<\eps<1/2$ there is $C_\eps$ with
$U_P(x)\le C_\eps x^{-1/2+\eps}$ for all sufficiently large $x$, then RH holds for every $L_k$ with
$h_k>0$. In particular $r=N$ detects every $L_k$ with $1\le k\le N$.
\end{theorem}

\begin{proof}
For $\Re s>1$, termwise integration and \eqref{hec:eq:DP} give
\[
 \int_1^\infty U_P(x)x^{-s}\,\dd x=\frac{D_P(s)}{s(s-1)}-\frac{h_0}{(s-1)^2}-\frac{g_P}{s-1},
\]
regular at $1$ by construction. At a zero $\rho$ in the open strip the residue is
$(h_0\ord_\rho\zeta_K+2\sum_kh_k\ord_\rho L_k)/(\rho(1-\rho))$, positive at a real zero detected by $P$
and nonzero at any detected nonreal zero; the Hecke $L$-functions have no zeros on $\Re s=1$. Apply
\cref{hec:lem:landau} to $g(x)=xU_P(x)$ with $a=1/2+\eps$, let $\eps\to0$, and use the functional
equations.
\end{proof}

The optimum determines the smallest pole term allowed by this class of detectors. The arithmetic estimate
concerns the angles of individual Gaussian prime powers, and is open already for $N=r=1$.

\subsection{Modular forms and symmetric powers}\label{hec:ss:sym}
{\emergencystretch=\textwidth\sources{grh-torsion-and-zero-data; grh_transfers_singular_moduli_and_symmetric_powers; grh_deficit_transfer_2026-11-25}}

The Landau step is not restricted to representations with finite image. For a holomorphic newform
without complex multiplication, the symmetric powers are cuspidal automorphic by Newton and Thorne
\cite{NewtonThorne2021Symmetric}, and a single positive prime sum tests a whole range of them. The
admissible weights are polynomials in the normalized Frobenius trace, and the size of the main term for a
given sensitivity is governed by extremal problems for the Sato--Tate measure.

\paragraph{Notation.}
Write $t=2\cos\theta$ and
\[
 C_j(t)=\frac{\sin((j+1)\theta)}{\sin\theta},\qquad
 \dd\mu(t)=\frac{\sqrt{4-t^2}}{2\pi}\,\dd t\quad(-2\le t\le2),
\]
with the continuous values at the endpoints. Thus $C_{-1}=0$, $C_0=1$, $tC_j=C_{j+1}+C_{j-1}$, the $C_j$
are monic and orthonormal for the Sato--Tate measure $\mu$, and
\begin{equation}\label{hec:eq:CG}
 C_iC_j=\sum_{k=0}^{\min(i,j)}C_{i+j-2k}.
\end{equation}
For a polynomial $P=\sum_ja_jC_j$ we write $[C_j]P=a_j$. The coefficient $a_0=\int P\,\dd\mu$ is the mean,
and $a_j$ is the weight given to the $j$th symmetric power. If $\alpha$ has absolute value one and
$t=\alpha+\alpha^{-1}$, then $C_j(t)=\sum_{k=0}^j\alpha^{j-2k}$ is the trace of the $j$th symmetric power
of $\operatorname{diag}(\alpha,\alpha^{-1})$, and $C_j(\alpha^m+\alpha^{-m})$ is the trace of its $m$th
power.

\paragraph{Three extremal problems.}
The first maximizes the weight of the standard $L$-function for a given mean; the other two minimize the
mean for a given weight on the highest symmetric powers.

\begin{proposition}[Largest trace coefficient; \stU]\label{hec:prop:extrE}
Let $d\ge1$. For a nonzero polynomial $f\ge0$ on $[-2,2]$ of degree at most $2d$, written
$f=\sum_{j=0}^{2d}a_jC_j$, the largest possible ratio $a_1/a_0$ is $u=2\cos(\pi/(d+2))$. The optimizer is
unique up to a positive scalar:
\[
 f_d(z)=\Bigl(\sum_{j=0}^dC_j(u)C_j(z)\Bigr)^2 ,
\]
and all its coefficients $a_0,\dots,a_{2d}$ are positive.
\end{proposition}

\begin{proof}
Let $\theta_i=i\pi/(d+2)$, $z_i=2\cos\theta_i$ and $w_i=2\sin^2\theta_i/(d+2)$ for $1\le i\le d+1$. The rule
$\int g\,\dd\mu=\sum_iw_ig(z_i)$ holds for polynomials of degree at most $2d+1$: divide $g$ by $C_{d+1}$,
which vanishes at all nodes; the quotient has degree at most $d$, so the integral of the multiple of
$C_{d+1}$ vanishes by orthogonality, and discrete orthogonality of sines gives the formula for the
remainder. Now $a_1/a_0=\int zf\,\dd\mu/\int f\,\dd\mu$ is a weighted average of the nodes, so it is at
most $z_1=u$, and equality forces a double zero of $f$ at every other node, which determines $f$ up to
scale. Discrete orthogonality gives $\sum_{j=0}^dC_j(z_1)C_j(z_i)=0$ for $i>1$, so $f_d$ attains equality.
Each $C_j(u)$, $0\le j\le d$, is positive, and by \eqref{hec:eq:CG} every $C_m$ with $0\le m\le2d$ occurs in
some product $C_jC_k$ with $j,k\le d$; hence $a_m>0$.
\end{proof}

For $d=2$ one has $u=\sqrt2$, $f_2(z)=(z^2+\sqrt2z)^2$, and the coefficients are
$(a_0,\dots,a_4)=(4,4\sqrt2,5,2\sqrt2,1)$.

For the other two problems fix $k\ge1$, let $u_k=(k+1)/k$, and for $0\le a\le u_k$ put
$Q_a=C_k+aC_{k-1}$. Its zeros $z_1<\dots<z_k$ are simple and lie in the open interval $(-2,2)$, except that
$z_1=-2$ when $a=u_k$: for $k\ge2$, at the zeros of $C_{k-1}$ the values $Q_a=C_k$ alternate in sign, and
$Q_a(2)=k+1+ak$, $(-1)^kQ_a(-2)=k+1-ak$ give the two outer roots; for $k=1$ the zero is $-a$. Let $\ell_i$
be the cardinal interpolation polynomials at the $z_i$ and $\lambda_i=\int\ell_i\,\dd\mu$. For
$\deg V\le2k-2$, division by $Q_a$ leaves a quotient of degree at most $k-2$, orthogonal to $Q_a$, and the
remainder $\sum_iV(z_i)\ell_i$. Hence
\begin{equation}\label{hec:eq:quadrature}
 \int V\,\dd\mu=\sum_i\lambda_iV(z_i)\quad(\deg V\le2k-2),\qquad\lambda_i=\int\ell_i^2\,\dd\mu>0 .
\end{equation}

\begin{proposition}[Least mean in odd degree; \stU]\label{hec:prop:odd}
Among monic polynomials of degree $2k-1$ that are nonnegative on $[-2,2]$, the least mean is $(k+1)/k$. The
unique minimizer is $R_k(t)=Q_{u_k}(t)^2/(t+2)$, and all its coefficients are positive:
\[
 [C_{2j}]R_k=\frac{k(k+1)-j(j+1)}{k^2}\quad(0\le j\le k-1),\qquad
 [C_{2j-1}]R_k=\frac{j^2}{k^2}\quad(1\le j\le k).
\]
\end{proposition}

\begin{proof}
Since $Q_{u_k}(-2)=0$, $R_k$ is a monic polynomial, nonnegative on the interval, with a simple zero at the
endpoint and $k-1$ double interior zeros. If $P$ is a competitor, $P-R_k$ has degree at most $2k-2$, and
\eqref{hec:eq:quadrature} with $a=u_k$ gives $\int(P-R_k)\,\dd\mu=\sum_i\lambda_iP(z_i)\ge0$. Equality
forces the endpoint zero and every interior double zero, which with monicity determine $R_k$. To check
the coefficients $r_j=[C_j]R_k$, multiply by $t+2$: the coefficient at $j$ becomes
$2r_j+r_{j-1}+r_{j+1}$, which is $1+u_k^2$ at even $j<2k$, $2u_k$ at odd $j<2k$ and $1$ at $2k$; these are
the coefficients of $Q_{u_k}^2$ by \eqref{hec:eq:CG}. In particular $r_0=u_k$.
\end{proof}

In even degree, minimizing the mean with only the top coefficient prescribed gives $C_k^2$, which omits
every odd symmetric power. Prescribing the coefficient of $C_{2k-1}$ measures the cost of including the
preceding power.

\begin{theorem}[Least mean in even degree; \stU]\label{hec:thm:even}
Let $P$ be monic of degree $2k$, nonnegative on $[-2,2]$, with $[C_{2k-1}]P=b\ge0$. The least possible mean
is
\[
 M_k(b)=\begin{cases}1+b^2/4,&0\le b\le2u_k,\\ 1-u_k^2+u_kb,&b\ge2u_k,\end{cases}
\]
and the unique minimizer is
\[
 P_{k,b}=\begin{cases}Q_{b/2}^2,&0\le b\le2u_k,\\ Q_{u_k}^2+(b-2u_k)R_k,&b\ge2u_k .\end{cases}
\]
For $b>0$ every coefficient $[C_j]P_{k,b}$ is positive.
\end{theorem}

\begin{proof}
For $0\le a\le u_k$ subtract $Q_a^2+(b-2a)C_{k-1}Q_a$ from $P$ to cancel the two highest terms. Applying
\eqref{hec:eq:quadrature} to the difference and using orthonormality gives the identity
\[
 \int P\,\dd\mu=1-a^2+ab+\sum_{i=1}^k\lambda_iP(z_i).
\]
Maximizing the quadratic lower bound over $[0,u_k]$ selects $a=\min(b/2,u_k)$ and gives $M_k(b)$. In the
first range $Q_a^2$ attains equality. For $b\ge2u_k$ the sum of $Q_{u_k}^2$ and $(b-2u_k)R_k$ is
nonnegative, has the prescribed two top coefficients and vanishes at all quadrature nodes. In the case of
interior nodes equality forces $k$ double zeros, hence the monic square. At $a=u_k$ equality forces
divisibility by $R_k$, and the remaining monic linear factor is determined by $b$. By \eqref{hec:eq:CG},
$[C_{2k}]Q_a^2=1$, all lower even coefficients equal $1+a^2$ and all odd coefficients equal $2a$;
\cref{hec:prop:odd} supplies the additional coefficients in the second range.
\end{proof}

For $0<b\le2u_k$ the same $M_k(b)$ is the sharp minimum when \emph{every} odd coefficient is required to be
at least $b$: a competitor has top odd coefficient $b'\ge b$ and $M_k(b')\ge M_k(b)$, while $Q_{b/2}^2$ has
all odd coefficients equal to $b$. Without a positive lower bound on some odd coefficient, the infimum of
the mean over polynomials with positive coefficients is $1$, approached as $b\downarrow0$ and not attained.
Under $P(2\cos\theta)\mapsto\sin\theta\,P(2\cos\theta)$ these are extremal problems for nonnegative sine
polynomials \cite{RogosinskiSzego1950Extremum}. The individual top-coefficient extremizers, including
\cref{hec:prop:odd} and the even-degree example $C_k^2$, are in \cite{DimitrovMerlo2002Nonnegative}, whose
next-to-top extremizer for odd sine degree is the case $b=2$ of \cref{hec:thm:even}; the case $k=1$ is
covered by \cite{RoysterSuffridge1970Typically}. The joint formula fixes both highest coefficients and
describes the transition at $b=2+2/k$.

\paragraph{The criterion for an elliptic curve.}
Fix an elliptic curve $E/\Q$ without complex multiplication and $d\ge1$. For a good prime $p$ let
$2\cos\theta_p=(p+1-\#E(\F_p))/\sqrt p$, $0\le\theta_p\le\pi$. Let $L_j(s)=L(s,\Sym^jE)$ in the unitary
normalization, with centre $1/2$, and $L_0=\zeta$. For $j\ge1$ these are cuspidal automorphic
$L$-functions \cite{NewtonThorne2021Symmetric}, entire with a functional equation and nonvanishing on
$\Re s=1$ \cite{JacquetShalika1976Nonvanishing}. With the extremizer $f=f_d=\sum_{j=0}^{2d}a_jC_j$ of
\cref{hec:prop:extrE} form the positive sum
\[
 S_f(x)=\sum_{\substack{p^k\le x,\ k\ge1\\p\ \mathrm{good}}}(\log p)\,f(2\cos(k\theta_p))
 \Bigl(\frac1{p^k}-\frac1x\Bigr),
\]
and, with $L_j^{(S)}$ the Euler product with the bad primes omitted,
\[
 D(s)=-\sum_{j=0}^{2d}a_j\frac{(L_j^{(S)})'}{L_j^{(S)}}(s),\qquad
 c_f=\lim_{s\to1}\Bigl(D(s)-\frac{a_0}{s-1}\Bigr)-a_0,\qquad
 R_f(x)=S_f(x)-a_0\log x-c_f .
\]
The bad Euler factors have normalized inverse roots of modulus at most one, so their removal introduces no
zeros in $\Re s>0$.

\begin{theorem}[Symmetric powers of an elliptic curve on fourth powers; \stR]\label{hec:thm:symE}
All nontrivial zeros of $L_0,\dots,L_{2d}$ have real part $1/2$ if and only if, for each $\eps>0$, there is
$B_\eps>0$ with
\[
 R_f(n^4)\le B_\eps n^{-2+4\eps}\qquad\text{for all sufficiently large integers }n .
\]
Under these equivalent conditions, $\abs{R_f(x)}\ll_{E,d,\eps}x^{-1/2+\eps}$ for all $x$.
\end{theorem}

\begin{proof}
\emph{Interpolation.} This is the chord interpolation of the first step of the proof of
\cref{ana:thm:sparse}. As a function of $t=1/x$, $S_f(1/t)$ is convex, each summand being a nonnegative
coefficient times $(p^{-k}-t)_+$. Between $t_-=(n+1)^{-4}$ and $t_+=n^{-4}$ it lies below its secant. The
secant of $-a_0\log t$ exceeds that function by at most $a_0(t_+-t_-)^2/(8t_-^2)=O_d(n^{-2})$, by the
interpolation remainder and the second derivative $a_0/t^2$. Since $x\asymp n^4$, the sampled bounds imply
$R_f(x)\le B'_\eps x^{-1/2+\eps}$ for all large $x$.

\emph{Zeros from the upper bound.} For $\Re s>1$ termwise integration gives
\begin{equation}\label{hec:eq:symmellin}
 \int_1^\infty R_f(x)x^{-s}\,\dd x=\frac{D(s)}{s(s-1)}-\frac{a_0}{(s-1)^2}-\frac{c_f}{s-1},
\end{equation}
because $C_j(2\cos(k\theta_p))$ is the trace of the $k$th power of the normalized Frobenius element on
$\Sym^j$. The definition of $c_f$ removes the singularity at $1$. At a zero $\rho$ in the open critical
strip the residue is
\[
 \frac{\sum_ja_j\ord_\rho L_j}{\rho(1-\rho)},
\]
which is positive at every real zero and nonzero at every nonreal zero, including a zero shared by several
factors and a central zero of any order. \Cref{hec:lem:landau} with $g(x)=xR_f(x)$ and $a=1/2+\eps$ shows
that the transform is holomorphic in $\Re s>1/2+\eps$; let $\eps\to0$ and use the functional equations.

\emph{The reverse implication.} Assume the zeros are on the line and fix $0<\eps<1/2$. Invert
$D(s)/(s(s-1))$ on $\Re s=2$ against $x^{s-1}$; the inverse of each prime-power term is
$(p^{-k}-x^{-1})_+$. Shift the line to $\Re s=1/2+\eps$. The double pole at $1$ contributes
$a_0\log x+c_f$. For the shift one uses the order-one Hadamard products of the fixed $L$-functions, their
gamma factors, and the local zero count: subtracting the logarithmic derivative at $2+it$ from that at
$\sigma+it$ gives an absolutely convergent sum over zeros, in which the zeros with $\abs{\Im\rho-t}\le1$
contribute $O_{E,d,\eps}(\log(2+\abs t))$ because $\abs{\sigma+it-\rho}\ge\eps$, and the more distant zeros
contribute the same by summing unit bands with weights $O((1+\abs{t-\Im\rho})^{-2})$. Thus
$D(\sigma+it)=O_{E,d,\eps}(\log(2+\abs t))$ uniformly on the strip for $\abs t\ge2$, the quadratic
denominator gives absolute convergence on the new line and vanishing horizontal integrals, and
$\abs{R_f(x)}\ll_{E,d,\eps}x^{-1/2+\eps}$.
\end{proof}

For $d=2$ the weights $(4,4\sqrt2,5,2\sqrt2,1)$ give an explicit target on the fourth powers for the first
four symmetric powers. For the curve $y^2+y=x^3-x$, of discriminant $37$ and nonintegral $j$-invariant
$48^3/37$, counting points at all good primes up to $10^4$ gives $S_f(16)=1.17616326$ and
$S_f(10^4)=19.38460159$ (\artifact{grh-torsion-and-zero-data/elliptic_samples.py}; floating-point sums).
{\emergencystretch=14em\par}

\begin{proposition}[Congruence weights for the elliptic test; \stU]\label{hec:prop:66}
Let $q\ge1$ and let $g$ be a real function on $(\Z/q\Z)^\times$. With $\pi_S(X)$ the number of good primes
$p\le X$ prime to $q$,
\[
 \lim_{X\to\infty}\frac1{\pi_S(X)}\sum_{\substack{p\le X\ \mathrm{good}\\p\nmid q}}
 \bigl(f_2(2\cos\theta_p)-g(p\bmod q)\bigr)^2
 =66+\frac1{\varphi(q)}\sum_{b\in(\Z/q\Z)^\times}(4-g(b))^2 .
\]
\end{proposition}

\begin{proof}
In each reduced residue class the average of $C_j(2\cos\theta_p)$ tends to zero for $1\le j\le8$. To see
this, expand the indicator of the class in Dirichlet characters and apply the prime number theorem to the
standard $L$-functions of $\Sym^jE$ twisted by these characters. Each twist is cuspidal of degree
$j+1>1$; its $L$-function is holomorphic and nonzero on $\Re s=1$, and its normalized prime-power traces
are bounded by $j+1$, which gives the prime number theorem by the usual Tauberian argument (add a large
multiple of $-\zeta'/\zeta$ to dominate the real and imaginary parts of the coefficients by nonnegative
ones, then subtract the ordinary prime number theorem). Prime powers of exponent at least two cost $o(X)$
in the logarithmically weighted sum, and partial summation removes the logarithm. Consequently the average
of any polynomial of degree at most $8$ in $2\cos\theta_p$ on a class is its Sato--Tate integral, and
orthonormality gives the squared fluctuation $(4\sqrt2)^2+5^2+(2\sqrt2)^2+1^2=66$.
\end{proof}

The proposition has the shape of \cref{nab:thm:periodic}, with the Sato--Tate measure in place of the Haar
measure of a finite Galois group; it is not an instance of that theorem, because the angles $\theta_p$ are
not Frobenius classes of a finite extension, and its proof uses the prime number theorem for the twisted
symmetric powers instead of the Chebotarev theorem.

\paragraph{The criterion for $\Delta$.}
Let $\Delta(z)=\sum_{n\ge1}\tau(n)e^{2\pi inz}$, and at each prime $p$ let
\[
 t_{p,m}=\alpha_p^m+\alpha_p^{-m},\qquad\alpha_p+\alpha_p^{-1}=\tau(p)p^{-11/2},\qquad\abs{\alpha_p}=1,
\]
the last by Deligne's bound \cite{Deligne1974Weil}. The level is one, so every Euler factor is unramified.
Put $L_j(s)=L(s,\Sym^j\Delta)$ in the unitary normalization and $L_0=\zeta$. Let $P=\sum_{j=0}^da_jC_j$ with
$a_j>0$ for $0\le j\le d$, and
\begin{align*}
 S_P(x)&=\sum_{p^m\le x}(\log p)P(t_{p,m})\Bigl(\frac1{p^m}-\frac1x\Bigr),\qquad
 D_P(s)=-\sum_{j=0}^da_j\frac{L_j'}{L_j}(s),\\
 c_P&=\lim_{s\to1}\Bigl(D_P(s)-\frac{a_0}{s-1}\Bigr)-a_0,\qquad
 R_P(x)=S_P(x)-a_0\log x-c_P .
\end{align*}
Then $D_P$ is the Dirichlet series with coefficient $(\log p)P(t_{p,m})$ at $p^m$. The extremal
polynomials $f_d$, $R_k$ and $P_{k,b}$ ($b>0$) above are admissible and are in addition pointwise
nonnegative.

\begin{theorem}[Symmetric powers of $\Delta$; \stR]\label{hec:thm:symDelta}
For each fixed $P$ as above the following are equivalent.
\begin{enumerate}[label=\textup{(\arabic*)}]
\item Every nontrivial zero of every $L_j$, $0\le j\le d$, lies on $\Re s=1/2$.
\item For every $0<\eps<1/2$ there are $K>0$ and $X\ge2$ such that $R_P(x)\le Kx^{-1/2+\eps}$ for
$x\ge X$.
\end{enumerate}
\end{theorem}

\begin{proof}
The Mellin identity \eqref{hec:eq:symmellin} holds with $D_P$, $a_0$ and $c_P$, and the residue at a zero
$\rho$ in the open strip is $\sum_ja_j\ord_\rho L_j/(\rho(1-\rho))$, nonzero at every such zero and positive
at a real one. The implication from (2) to (1) is therefore the second step of the proof of
\cref{hec:thm:symE}. Conversely, let $\psi_P(y)=\sum_{p^m\le y}\log p\,P(t_{p,m})$. Under (1) the explicit
formula for the fixed symmetric-power $L$-functions gives $\psi_P(y)=a_0y+O_P(\sqrt y\log^2(2y))$
\cite{RouseThorner2017Explicit}. Since $S_P(x)=\int_1^x\psi_P(y)y^{-2}\,\dd y$, the integral of
$(\psi_P(y)-a_0y)y^{-2}$ over $[1,\infty)$ converges, its tail is $O_P(x^{-1/2}\log^2(2x))$, and its value
is $c_P$, as one sees by expanding $D_P(s)=s\int_1^\infty\psi_P(y)y^{-s-1}\,\dd y$ at $1$.
\end{proof}

\begin{proposition}[Prime-only form; \stU]\label{hec:prop:primeonlyDelta}
Let $R_P^{(1)}(x)$ be defined as $R_P(x)$ with only the terms $m=1$ in $S_P(x)$ and with $c_P$ replaced by
$c_P^{(1)}=c_P-\sum_{p,\,m\ge2}(\log p)p^{-m}P(t_{p,m})$. Then
\[
 R_P^{(1)}(x)-R_P(x)=\frac{2B_P+o(1)}{\sqrt x},\qquad B_P=\sum_{j=0}^d(-1)^ja_j .
\]
For the odd-degree extremizer $R_r$ of degree $2r-1$ one has $B_P=(r+1)(2r+1)/(6r)$. Consequently the
criterion of \cref{hec:thm:symDelta} can be stated with primes alone.
\end{proposition}

\begin{proof}
The correction to the constant converges absolutely. Put
\[
 H(y)=\sum_{p^m\le y,\ m\ge2}\log p\,P(t_{p,m}).
\]
The terms with $m\ge3$ are $O_P(y^{1/3}\log^2(2y))$. For
the squares, $t_{p,2}=t_{p,1}^2-2$, and the Sato--Tate theorem for $\Delta$ with partial summation gives
$H(y)=B_P\sqrt y+o(\sqrt y)$ with $B_P=\int P(t^2-2)\,\dd\mu(t)$. Since
\[
 C_j(t^2-2)=\sum_{k=0}^j(-1)^kC_{2(j-k)}(t),
\]
only the last term has nonzero mean, which gives the formula for $B_P$; inserting the coefficients of
\cref{hec:prop:odd} gives the value for $R_r$. Partial summation gives
\[
 R_P^{(1)}(x)-R_P(x)=\frac{H(x)}x+\int_{(x,\infty)}\frac{\dd H(y)}y=\int_x^\infty\frac{H(y)}{y^2}\,\dd y .
\]
\end{proof}

For $d=1$ the weight is $P(t)=2+t$, and the criterion tests the zeros of the degree-two $L$-function of
$\Delta$ together with those of $\zeta$.

\paragraph{A nonnegative deficit weight for a single $L$-function.}
A different normalization isolates one $L$-function with nonnegative weights and without a centring
constant. For $c>1$ let
\[
 K_c(y)=\begin{cases}1-c^{-1},&0<y\le1,\\ y^{-1}-c^{-1},&1<y\le c,\\0,&y>c,\end{cases}
\]
so that for a coefficient sequence $a(n)$ with $T(x)=\sum_{n\le x}a(n)(1/n-1/x)$ one has
$T(cx)-T(x)=x^{-1}\sum_{n\le cx}a(n)K_c(n/x)$. The transform of this finite difference carries the factor
$H_c(s)=(c^{s-1}-1)/(s(s-1))$, which is positive on the positive real axis and has no zeros in
$0<\Re s<1$; the one-sided finite-difference detector built on it is \cref{ana:thm:fd}.

\begin{theorem}[Deficit detector; \stR]\label{hec:thm:deficit}
Let $L(s,\pi)$ be a normalized self-dual cuspidal automorphic $L$-function of degree $d$ over $\Q$ whose
Euler traces $t(p^k)$ are real and bounded by $d$ in absolute value, with the usual continuation,
functional equation, order-one completion and nonvanishing on $\Re s\ge1$. Put $w(p^k)=d-t(p^k)\ge0$.
Assume \cref{src:thm:dgrh}. Then RH for $L(s,\pi)$ is equivalent to the statement that for every $\eps>0$
\begin{equation}\label{hec:eq:deficit}
 \frac1x\sum_{p\le2x}(d-t(p))\log p\;K_2(p/x)-d\log2\le C_{\pi,\eps}x^{-1/2+\eps}
\end{equation}
for all sufficiently large $x$. The same holds with the pointwise bound replaced by the corresponding
bound for the mean of the positive part over $[X,2X]$.
\end{theorem}

\begin{proof}
Up to finitely many harmless Euler corrections, the Dirichlet series of $w(p^k)\log p$ is
\[
 A(s)=-d\,\frac{\zeta'(s)}{\zeta(s)}+\frac{L'(s,\pi)}{L(s,\pi)} .
\]
It has residue $d$ at $1$ and, since $\zeta$ is zero-free in $\Re s>1/2$, a simple pole with positive
residue, the multiplicity, at each zero of $L(s,\pi)$ to the right of the critical line. The
finite-difference detector of \cref{ana:thm:fd} with $\lambda=2$, applied to the coefficients $-w$ (the
opposite orientation, \cref{ana:rem:sign}), therefore gives the absence of such zeros from
\eqref{hec:eq:deficit}, or from its averaged form, and the functional equation gives RH. The converse
follows from the contour shift for $T(2x)-T(x)$, whose multiplier $H_2$ is $O((1+\abs t)^{-2})$ on vertical
strips. The prime powers of exponent at least two contribute $O(x^{-1/2}\log^2(2x))$ to the left side of
\eqref{hec:eq:deficit}, so passing from prime powers to primes changes neither implication.
\end{proof}

For $\Delta$ one has $d=2$ and $t(p)=\tau(p)/p^{11/2}$. For the $r$th symmetric power, $d=r+1$ and
$t_r(p)=C_r(t(p))$, so the weight $r+1-C_r(\tau(p)p^{-11/2})$ gives a nonnegative target for RH of
$L(s,\Sym^r\Delta)$ alone. In contrast with \cref{hec:thm:symDelta}, the tested function enters with a
negative coefficient, the test detects a single symmetric power, and the residue at $1$ is the full degree.
Neither criterion bears on the order of vanishing at the centre.

\subsection{Open questions}\label{hec:ss:open}
{\emergencystretch=\textwidth\sources{grh_transfers; grh-lean-transfers; grh_dihedral_47_71; grh_transfers_density_gaps_and_zero_tests; mcnew_density_and_grh_transfers; mcnew_counterexample_and_grh_transfers; grh-methods-polynomial-density; grh-torsion-and-zero-data; grh_transfers_singular_moduli_and_symmetric_powers; mcnew-density-and-dirichlet-grh; grh_transfers_lubin_tate_and_graph_towers}}

Every reduction in this section leaves a one-sided prime estimate unproved. The questions below ask for
these estimates in their smallest instances, and for the extensions of the unconditional results that the
proofs suggest.

\begin{enumerate}[label=\textup{(\arabic*)}]
\item \emph{The conductor-$23$ estimate.} Can a covolume comparison of the kind used in the source proof,
built over the CM base $\Q(\sqrt{-23})$ with data indexed by its ideal classes, produce any one of the
equivalent one-sided estimates of \cref{hec:ss:cubic}: the bound \eqref{hec:eq:principalbound} for the
principal-class remainder, the eventual sign $\mathcal T(x_j)\le0$ of \cref{hec:thm:sign}, or the
irreducible-prime bound \eqref{hec:eq:irred23}? The construction would need integral forms whose local
divisibility distinguishes the identity from a three-cycle inside the same quadratic splitting condition
(the principal form $u^2+uv+6v^2$ represents the completely split primes and $2u^2+uv+3v^2$ the primes at
which the cubic is irreducible), enough independent such forms, a determinant lower bound with controlled
denominators, and a matching archimedean main term. \Cref{hec:prop:modulus,hec:prop:meansq} show that
this information cannot come from a fixed set of residue classes. The same question is open for the weight
$r_p-1$ of the field $\Q(\sqrt[3]2)$ in \cref{hec:thm:cuberoot}.
\item \emph{A selective event.} Can the comparison be realized for the binary event of
\cref{hec:thm:twistedcubic}, with its negative coefficient on $\rho\eta$ and its vanishing coefficient on
$\rho$?
\item \emph{Other sign tests.} For which other self-dual Artin or modular $L$-functions of degree two does
the bias from prime squares exceed the total contribution of the zeros, as in \eqref{hec:eq:limsup}, so
that the Riemann hypothesis for the function becomes equivalent to an eventual sign of a prime sum? The
curvature bound of \cref{zer:prop:eta} gives a verifiable sufficient condition for positive eta products.
\item \emph{Conductor $47$.} Can the prime weight of the class of $Q_1$ in \cref{hec:thm:primes47} be
realized as a local divisibility contribution of an integral theta lift, with the determinant lower bound
retained? Reproducing its mean $1/5$ gives no control of its error.
\item \emph{Prime dihedral degrees.} For a dihedral Hilbert class field of prime degree $\ell>7$ over its
quadratic subfield, a single rotation class still prevents cancellation of zero multiplicities. For which
explicit fields can the real-axis hypothesis be proved by positivity of the Mellin kernel of every
constituent, as in \cref{hec:prop:pos71}?
\item \emph{Real zeros of pencils.} Which pairs of newforms of equal conductor admit a quotient with a
single simple pole on a Fricke quotient, as in the proof of \cref{hec:thm:real47}? When the modular
function has degree greater than one, can the number of real zeros of a linear combination be bounded by
the number of turns of the ratio of the kernels?
\item \emph{Cubic Jensen inequalities.} Which property of the conductor-$23$ kernel yields the cubic
Jensen inequalities at every shift without locating zeros? By
\cref{hec:thm:finitejensen,hec:thm:stride,hec:thm:sexticlimit} no argument that uses only the sign of
$(\log K(\sqrt x))''$ or of the fourth derivative of $\log K$ can do so, and an argument uniform in the
degree would be needed for the full critical-line statement.
\item \emph{The sextic family.} For each $c>0$, what is the least degree $d$ for which $Q_{d,c}$ of
\cref{hec:thm:sexticlimit} fails to be real-rooted? The bound $d>2+8/(9c^2)$ is sufficient.
\item \emph{Gaussian angles.} Can the local divisibility calculation in a theta lift be made sensitive to
one nonconstant angular harmonic of Gaussian primes, that is, can the bound of \cref{hec:thm:angular} be
proved for $N=r=1$?
\item \emph{Traces of Frobenius.} Can a height or divisibility inequality give the upper bound of
\cref{hec:thm:symDelta} for the degree-one weight $2+\tau(p)p^{-11/2}$, or the bound of \cref{hec:thm:symE}
for the curve $y^2+y=x^3-x$ and $d=1$? \Cref{hec:prop:66} identifies the part of the local data that no
fixed modulus sees.
\item \emph{Central nonvanishing.} For quadratic characters, can sampling sites depending on the
conductor, with a number of samples $k=k(q)$ tending to infinity and uniform estimates for the logarithmic
derivatives at these sites, force a central multiplicity bound below two? By \cref{hec:thm:kernels} no
bounded number of samples can.
\end{enumerate}

\emergencystretch=\hecsavedemergencystretch\relax
